% concatenated from limitations_of_qaoa.tex
\documentclass[a4paper,onecolumn,11pt,unpublished]{quantumarticle}
\pdfoutput=1
\usepackage[utf8]{inputenc}
\usepackage[english]{babel}
\usepackage[T1]{fontenc}
\usepackage{amsmath}
\usepackage{amssymb}
\usepackage{physics}

\usepackage{hyperref}
\usepackage[table]{xcolor}
\usepackage{colortbl}
\usepackage{booktabs}
\usepackage{csquotes}
\usepackage{enumerate}
\usepackage{enumitem}
% \usepackage{subfigure}
\usepackage{amsthm}
% in the very top of your preamble, after \documentclass
\usepackage{caption}
\usepackage{subcaption}

\usepackage{selinput}
\SelectInputMappings{
  adieresis={ä}, % etc.
  acute={´},
}
\usepackage{amsmath}  % if not already loaded
\newcommand{\Unif}{\mathcal{U}}
\usepackage[english]{babel}
\usepackage[T1]{fontenc}
\usepackage{lipsum}
\usepackage{mdframed}
\usepackage{booktabs}
\usepackage{adjustbox}
\usepackage{longtable}

% \usepackage[utf8]{inputenc}
% \usepackage[T1]{fontenc}
\usepackage[
  backend=biber,
  style=numeric,
  sorting=none      % or 'nyt' if you want author–year sorting
]{biblatex}
\addbibresource{ref.bib}
% 1) load xcolor & define the exact violet used by Quantum
\usepackage{xcolor}
\usepackage[table]{xcolor}
\usepackage{booktabs}
\usepackage{listings}
% in your document preamble:
\usepackage{siunitx}
% \usetikzlibrary{positioning}
\usepackage{multicol}
% allow TeX to stretch lines a bit more before complaining
\setlength{\emergencystretch}{17em}



% define a style for appendix Python
\lstdefinestyle{TinyPython}{
  language        = Python,
  basicstyle      = \ttfamily\fontsize{5pt}{6pt}\selectfont,  % 5 pt type
  numbers         = left,
  numberstyle     = \ttfamily\fontsize{5pt}{6pt}\selectfont,
  stepnumber      = 1,
  numbersep       = 3pt,
  breaklines      = true,
  breakatwhitespace = true,
  showstringspaces = false,
  frame           = single,
  xleftmargin     = 0pt,
  xrightmargin    = 0pt,
  aboveskip       = 0pt,
  belowskip       = 0pt,
  tabsize         = 2,
  columns         = flexible,
  keepspaces      = true,
}
%─────


% 3) (optional) colour all section‐heads
% \usepackage{sectsty}
% \sectionfont{\color{quantumviolet}}     % sections
% \subsectionfont{\color{quantumviolet}}  % subsections

\usepackage{pifont}  
\usepackage{newunicodechar}
\newunicodechar{✖}{\ding{55}}
\usepackage[margin=1in]{geometry}   % Adjust page margins
\usepackage{setspace}              % For line spacing
\usepackage{amssymb} 
\usepackage{bm}
\usepackage{amsmath}
\usepackage{amsthm}  % Math symbols              % Theorem and proof environments
\usepackage{tikz}            % TikZ is required by quanti
\usepackage{quantikz} 
\usepackage{enumitem}% For drawing quantum circuits
\usepackage{setspace}              % For line spacing
\usepackage{bm}
\usepackage{amsmath,amssymb,amsthm}             % Theorem and proof environments
\usepackage{tikz}            % TikZ is required by quanti
\usepackage{quantikz} 
\usepackage{enumitem}% For drawing quantum circuits
\usepackage{graphicx} 
\usepackage{algpseudocode}
\usepackage{algorithm}
\usepackage{algorithmicx}
% \usepackage{lmodern}
% --- Math & theorem environments ---
% \usepackage{amsmath,amssymb}
\usepackage{amsthm}
\usepackage{xcolor}
% tcolorbox with “most” libraries for styling
\usepackage[most]{tcolorbox}
\usepackage{pgfplots}
\pgfplotsset{compat=1.17}
\newcommand{\poly}[1]{\mathrm{poly}\bigl(#1\bigr)}
% \newcommand{\poly}[1]{\mathrm{poly}\bigl(#1\bigr)}
% enumitem so that leftmargin=* works in enumerate

\usepackage{booktabs}
% Nicer table rules
\usepackage{float}                 % For [H] float placement
\usepackage{hyperref}              % Clickable links/references
\usepackage{url}
\usepackage{xcolor}                % For colored text (optional)
\usepackage{lscape}                % For landscape pages if needed
\usepackage{caption}
\usepackage{subcaption}
\usepackage{etoolbox}
% patch the abstract environment to be ragged-right
\AtBeginEnvironment{abstract}{\raggedright}
\usepackage{amsmath,amssymb,amsthm}
\newtheorem{theorem}{Theorem}
\newtheorem{lemma}[theorem]{Lemma}
\newtheorem{proposition}[theorem]{Proposition}
\newtheorem{corollary}[theorem]{Corollary}
\newtheorem{definition}[theorem]{Definition}

\theoremstyle{remark}
\newtheorem{remark}[theorem]{Remark}


% Macros
\newcommand{\OH}{\mathcal{H}_{\mathrm{OH}}}
% \newcommand{\ket}[1]{\left|#1\right\rangle}
% \newcommand{\bra}[1]{\left\langle#1\right|}
% \newcommand{\braket}[2]{\left\langle #1\,\middle|\,#2 \right\rangle}
% \newcommand{\ketbra}[2]{\left|#1\right\rangle\!\left\langle#2\right|}
\newcommand{\spec}{\mathrm{spec}}

% ---------- REDUCE-QAOA VERTICAL SPACE-QAOA BETWEEN EQUATIONS ----------
\makeatletter
\g@addto@macro\normalsize{%
    \abovedisplayskip 3pt plus 1pt minus 1pt%
    \abovedisplayshortskip 3pt plus 1pt minus 1pt%
    \belowdisplayskip 3pt plus 1pt minus 1pt%
    \belowdisplayshortskip 3pt plus 1pt minus 1pt%
}
\makeatother

\def\BibTeX{{\rm B\kern-.05em{\sc i\kern-.025em b}\kern-.08em
    T\kern-.1667em\lower.7ex\hbox{E}\kern-.125emX}}

% % before the \maketitle:
% \pretitle{%
%   \centering
%   \color{quantumviolet}\Huge\bfseries
% }
% \posttitle{\par\vspace{1em}}

% \preauthor{\centering\large}

\author{Chinonso Onah}
\affiliation{Volkswagen AG, Berliner Ring 2, 38440 Wolfsburg, Germany}
\affiliation{Department of Physics, RWTH Aachen University, 52056 Aachen, Germany}
\email{chinonso.calistus.onah@volkswagen.de}
% \email{c.onah@extern.fz-juelich.de}

\author{Kristel Michielsen}
\affiliation{Forschungszentrum Jülich, Germany}
\affiliation{Universit\"at zu K\"oln, 50923 K\"oln, Germany}

\title{Fundamental Limitations of QAOA on Constrained Problems and a Route  to Exponential Enhancement}



\begin{document}
\maketitle
% \title{}   % we won’t actually call \maketitle




\begin{abstract}
We study fundamental limitations of the generic Quantum Approximate Optimization Algorithm (QAOA) on constrained problems where valid solutions form a low-dimensional manifold inside the Boolean hypercube, and we present a provable route to exponential enhancements via constraint embedding. Focusing on permutation-constrained objectives, we show that the standard “generic’’ QAOA ansatz faces an intrinsic feasibility bottleneck: even after optimizing angles, the feasible probability mass of depth--\(p\) circuits remains controlled by the hypercube embedding scale
\(|\Pi|/2^N\), up to the overlap-corrected shallow-depth amplification
governed by the light-cone overlap number \(\chi_p\). In the one-dimensional nearest-row geometry, this includes all \(p=o(n/\log n)\). Our analysis triangulates this separation via four complementary bounds on generic QAOA: (i) Walsh–Fourier/Krawtchouk estimates showing that permutation one–hot constraints have exponentially small low-degree and high-degree Fourier mass; (ii) a cost-angle averaging argument that pins the typical-angle feasible baseline at $|\Pi|/2^N$; (iii) fourth-moment ($L_4$) bounds that control typical fluctuations around this baseline; and (iv) light-cone locality arguments that show shallow circuits cannot build the global correlations demanded by permutation constraints. Against this envelope we place a minimal constraint-enhanced kernel (CE–QAOA) that operates directly inside a product one–hot subspace and mixes with a block–local $XY$ Hamiltonian.  For permutation constrained problems on \(n\) variables, we prove an
angle-robust exponential enhancement of the form
\begin{equation}
\label{eq:expp}
  \frac{P_{1}^{\mathrm{(CE)}}}{P_{1}^{\mathrm{(gen)}}}
  \;\ge\;
  \exp\!\bigl[\Theta(n^2)\bigr]
\end{equation}
at \(p=1\). Beyond one layer and throughout overlap-controlled sublinear-depth windows \(  p=o\!\left(\frac{n}{\ln n}\right),\) the general feasibility ratio is modified by the overlap number \(\chi_p\) as $\exp\!\bigl[\Theta(\frac{n^2}{\chi_p})\bigr]$.% to give the corresponding overlap-corrected light-cone exponent.
\end{abstract}





\section{Introduction}
\label{sec:introduction}

Combinatorial optimization problems (COPs) with global constraints pose a specific challenge for shallow variational quantum circuits. Feasible solutions occupy a vanishing fraction of the full Hilbert space, so any method that explores the entire hypercube pays an intrinsic feasibility penalty. This paper isolates that bottleneck for the Quantum Approximate Optimization Algorithm (QAOA)\cite{Farhi2014QAOA} and develops a transparent way to assess the state prepared by the algorithm for arbitrary shallow schedules of length $p = O(n)$. QAOA is arguably one of the most studied Quantum algorithms in recent years and numerous variants have emerged. In this paper, we shall refer to the original construction in Ref. \cite{Farhi2014QAOA} as  the generic QAOA ansatz in order to facilitate a clear comparison with a minimal constraint-aware variant that bakes constraint structure and symmetries into its quantum dynamics. The distinction is made explicit in Subsection \ref{sec:back}.

\medskip
\noindent
Instead of monitoring only expectation values and approximation ratios, we study the full probability profile of depth-\(p\) QAOA states through four lines of analyses. First, a Walsh–Fourier calculation on the Boolean cube reveals that the one-hot constraints behind permutation encodings have exponentially small low-degree Fourier mass. At \(p=1\) the \(X\)-mixer acts by pure phases in this Walsh/Fourier basis, so the generic ansatz cannot significantly lift feasible mass above the uniform baseline, up to polynomial slack, and this holds uniformly over angles. Second, a gentle cost-angle averaging argument, applicable after a simple rescaling that holds for standard integer or rational instances, pins the feasible baseline at \(|\Pi|/2^N\). Third, a hypercontractive ($l_4$) bound on the output amplitudes then show that most angles concentrate around this baseline, which is useful for understanding shallow-depth training without assuming anything about gradients. Finally, locality enters through light-cone growth. With commuting, \(r\)-local diagonal costs, information propagates at most one neighborhood per layer. A dependency–graph argument turns this into a bound on the joint probability that all row and column constraints hold, and the same limitation persists beyond one layer in any depth window where the overlap-corrected light-cone exponent remains larger than the
encoded-manifold baseline cost \(n\ln n\).  In particular, for the
one-dimensional nearest-row geometry this includes \(
  p=o\!\left(\frac{n}{\ln n}\right).\) The resulting picture is that the generic ansatz faces a structural ceiling that does not disappear at depth one and continues to matter at larger but still shallow depths.

\medskip
\noindent
Against this envelope, we place the recently introduced variant of QAOA called the Constraint-Enhanced Quantum Approximate Optimization algorithm (CE–QAOA)\cite{onahce}. This is a constraint-enhanced kernel that works directly on the one-hot manifold with explicit symmetry enhancements derived from problem structure. The ingredients include an ancilla-free \(W\)-state initialization, a normalized block–\(XY\) mixer that preserves the encoded space and penalty derived co-design symmetry. This framework provides an existence result needed to benchmark the probability profile of generic QAOA. Thus, we take the bounds derived for CE–QAOA kernel from prior work~\cite{onahce} as a minimal exemplar of problem–algorithm co-design as a route to provable exponential enhancement. % Comparing the CE-QAOA lower envelope with the  upper envelope from generic QAOA yields an angle-robust and depth-matched exponential enhancement in feasible mass.  that paints a clear picture of the path to exponential enhancements through problem-algorithm co design. It provides, for any basis vector $\mid x^\star\rangle$,  the bound
% \(
%   \bigl|\langle x^{\star}\!\mid U_{M}(\boldsymbol\beta^{\star})\,U_{C}(\gamma^{\star})\,\psi_{0}\rangle\bigr|^{2}
%   \;\ge\; \frac{c}{n^{n}}, c>0,
% \)
% at p=1. This supplies a concrete feasibility guarantee %

\medskip
\noindent
For any permutation constrained problem of size $n$, its one-hot encoding requires \(N=n^2\) qubits (Def \ref{def:kernel-requirement}). The resulting enhancement is quantified in terms of $\exp(\Theta(n^2))$ separation between the constraint enhanced variant and the generic QAOA in Eq.~\eqref{eq:expp} and proved later in Theorem \ref{thm:p1-sep} for $p=1$ and extended in Theorem~\ref{thm:generic-const} to
overlap-controlled sublinear depths respectively. It is worth noting that we have only used the minimal, most conservative guarantee from Ref. \cite{onahce}. Our aim is to place the minimal guarantee from the constraint-enhanced kernel against a sharp upper envelope for the generic ansatz and to show how even a modest in-subspace design changes the feasibility profile at shallow depth. %In this sense, the kernel functions as a probe of the envelope rather than a new construction.

% \medskip
% \noindent
% In the spirit of parameter transfer for QAOA, we prove that angles which are optimal for the \emph{generic} ansatz at depth \(p{=}1\) already guarantee an exponential boost when reused verbatim inside the CE–QAOA kernel. Writing
% \(
% P_{\max}^{\textup{(gen)}}=\max_{\beta,\gamma}P_{\Pi}^{\textup{(gen)}}(\beta,\gamma)
% \) our analysis shows that for every \(n\ge2\) and every $(\gamma, \beta)$
% pair \[
% P^{\textup{(CE)}}
% \;\ge\;
% \sqrt{2\pi n}\,e^{\,n}\;P_{\max}^{\textup{(gen)}}.
% \]
% This “parameter-transfer amplification’’ complements our envelope bounds and explains the robust gains observed in the noiseless circuit comparisons reported in Figures~\ref{fig:wi4_counts}. 


% Our assumptions are mild. We require only that \(H_C\) is diagonal, commuting, \(r\)-local, and of bounded degree, and that the mixer is the standard transverse field. The typical-angle statements rely on a modest spectral rescaling that holds in common settings. On the CE–QAOA side we use the basic existence result which fits assignment and TSP-style encodings and follows from a simple problem–algorithm co-design~\cite{onahce}.

\medskip
\noindent
The rest of the paper proceeds as follows. Section~\ref{sec:related} places our work in context of recent research works in fundamental limitations of shallow Variational Quantum Algorithms (VQA). Section~\ref{sec:back} reviews the generic ansatz, recalls the definition of the CE–QAOA kernel and records their symmetry, mixing properties and the core existence result. In Section~ \ref{sec:main} we collect our main results. In Section~\ref{sec:methods} we discuss the materials and methods necessary to establish the angle-averaging, fourth-moment, and harmonic-analysis bounds that define the generic envelope. Section~\ref{app:lightcone} develops light-cone barriers for the generic QAOA ansatz at $p=1$ and Section~\ref{app:fixed_p} extends the light-cone barrier analyses to constant/sublinear depths. Section \ref{sec:conclusion} concludes our findings. 

%Supplemntary details are relgated to the Appendix~\ref{app:param-transfer}.






\subsection{Related Work}
\label{sec:related}

Variational quantum algorithms (VQAs), and QAOA in particular, have been studied along three complementary axes: (i) \emph{trainability and landscape structure}, (ii) \emph{locality/light-cone limitations at shallow depth}, and (iii) \emph{architecture- or problem-aware ans\"atze} that encode constraints or symmetries from the outset. 
On the landscape side, early work identified \emph{barren plateaus} as a key obstruction to gradient-based training in large systems~\cite{McClean2018BarrenPlateaus,Cerezo2021VQAReview}. 
More recent analyses quantify landscape \emph{information content}, revealing how much task-relevant signal survives in variational loss surfaces~\cite{PerezSalinasWangBonetMonroig2024}, and study the \emph{energy landscape} structure of small-depth QAOA instances~\cite{choy2024energylandscapesquantumapproximate}. 
The QAOA design space has also broadened beyond the original transverse-field mixer to \emph{alternating operator ans\"atze} tailored to problem structure~\cite{HadfieldHoggRieffel2022}, while symmetry constraints can both obstruct and guide expressivity~\cite{Bravyi_2020}. 
Mitigation strategies such as \emph{classical shadows} for gradient estimation~\cite{SackMedinaMichailidisKuengSerbyn2022} and \emph{warm starts}~\cite{Egger2021Warm} aim to stabilize training or inject structure into initial states.




A second line of work formalizes \emph{locality} as the principal barrier at small depths since shallow circuits can only build correlations within a \emph{light cone} whose radius grows at most linearly with depth. These are essentially Lieb--Robinson-type locality bounds~\cite{Hastings2010Locality}, and they underlie modern limitations for low-depth QAOA on sparse graphs. For bounded-degree graph families, Farhi, Gamarnik, and Gutmann proved that small-$p$ QAOA ``needs to see the whole graph,'' producing typical- and worst-case examples where performance is capped by local neighborhoods~\cite{FarhiGamarnikGutmann2020WholeGraph}. At constant levels $p=O(1)$, quantitative barriers have been established on sparse hypergraphs and spin-glass models~\cite{Basso2022}, while in a broader framework of \emph{local quantum algorithms}, limitations propagate up to $p<\varepsilon\log n$ for certain CSPs with the (coupled) overlap-gap property~\cite{GamarnikZadik2022QAOAOGP,Anschuetz_2024}. 
Orthogonally, results on concentration for shallow circuits show output observables of low-depth evolutions are tightly concentrated, limiting global coordination at small $p$~\cite{BravyiGossetKonig2018}. Collectively, these works support the heuristic that constant or slowly growing depth cannot generate the long-range correlations required by \emph{global} constraints.


A third, complementary, line of work focuses on \emph{architecture- and problem-aware}
ans\"atze that encode constraints, symmetries, or known structure directly at the
circuit level. In the QAOA setting, the alternating-operator framework extends the
original transverse-field mixer by replacing it with constraint-preserving unitaries
tailored to specific combinatorial structure~\cite{Hadfield2019AOA,HadfieldHoggRieffel2022}.
Constrained mixers and symmetry-projected variational families have been designed to
restrict the evolution to feasible manifolds or fixed symmetry sectors, thereby
reducing the effective search space and sometimes alleviating barren-plateau effects
~\cite{BaertschiEidenbenz2020,Fuchs2022ConstrainedMixers,Bravyi_2020}. Domain-specific encodings and mixers for routing and scheduling problems further illustrate how problem-aware design can unlock strong empirical performance on structured instances~\cite{Xie2024CVRP,tsvelikhovskiy2024symmetries}.


The current paper connects these three perspectives under a single quantitative
lens by focusing on the \emph{feasible probability mass} generated by shallow
generic versus problem-aware variational circuits under global combinatorial
constraints. Rather than bounding approximation ratios instance by instance,
we derive angle-uniform upper bounds on the total probability that a shallow,
\emph{generic} QAOA state satisfies all permutation-style one-hot constraints, for depths $p$ up to a sublinear fraction of $n$ on interaction hypergraphs whose incidence growth is polynomial in $n$. Against this envelope, we place a minimal \emph{constraint-enhanced} kernel (CE--QAOA) that encodes the instance in a constrained one-hot manifold and mixes \emph{within} it, thereby achieving a
depth-matched, \emph{angle-robust} exponential separation in feasible mass,
cf.\ Eq.~\eqref{eq:expp}, in a way that complements and connects prior
landscape, locality, and architecture-aware analyses. By construction, the
explicit problem–algorithm co-design in the kernel definition allows these
results to apply to a broad class of constrained problems satisfying
Definition~\ref{def:kernel-requirement}.





\subsection{Background}
\label{sec:back}

\subsubsection{Generic QAOA}
\label{subsec:generic-qaoa}
The generic QAOA~\cite{Farhi2014QAOA} prepares the state
\(
\ket{\psi_p(\vec\gamma,\vec\beta)}
=\bigl[\prod_{\ell=1}^{p}e^{-i\beta_\ell \sum_j X_j}\,e^{-i\gamma_\ell H_C}\bigr]\ket{+}^{\otimes N},
\)
i.e., a depth–$p$ alternation of a transverse–field $X$ mixer and a diagonal cost unitary on the computational basis. Parameters \((\vec\gamma,\vec\beta)\) are optimized (classically) to maximize an empirical objective. For the algorithm to be performant,  typical schedules need to be short (\(p=O(1)\) or \(p\ll N\)) without going to the adiabatic limit\cite{Lucas2014}. Common ideas to improve algorithmic performance include linear/adiabatic ramps\cite{montanezbarrera2024universalqaoa}, parameter transfer across sizes/instances to reduce training cost while keeping the same mixer/cost structures~\cite{Monta_ez_Barrera_2025T}, recursive elimination to iteratively fix variables based on low–depth runs, shrinking the instance before a final classical solve~\cite{Bravyi_2020, BaeLee2024RecursiveQAOA},  warmstarting QAOA to help landscapes and convergence~\cite{Egger2021Warm, carmo2025warmstarting}. 



\subsubsection{Constraint–Enhanced QAOA}
\label{sec:ce-qaoa}

The \emph{alternating–operator} (``QAOA+'') paradigm augments the generic X–mixer
by feasibility–preserving mixers that act invariantly on constraint projectors,
thereby confining evolution to structured subspaces
\cite{Hadfield2019AOA,Fuchs2022ConstrainedMixers}. The
\emph{Constraint–Enhanced QAOA} (CE–QAOA) introduced in Ref. \cite{onahce} follows this direction but makes the problem–algorithm co–design explicit. As can be seen in Definition~\ref{def:kernel-requirement}, the kernel specifies (i) the valid problem classes with the admissible penalty structure and symmetries, the mixer and initial
state.

% We introduce a kernel designed to operate on block one–hot manifolds and \emph{fix} the mixer family and initial states to match the encoding. 


\begin{definition}[CE--QAOA kernel]
\label{def:kernel-requirement}
An optimization instance $I$ belongs to the \emph{CE--QAOA kernel} if there exist
integers $n,m\in\mathbb N$ and the \emph{one-hot} encoder $\mathsf E_{\mathrm{1hot}}$
that initializes the dynamics in the constrained fixed–Hamming–weight space
\[
\OH \;=\; (\mathcal H_1)^{\otimes m},
\qquad
\mathcal H_1 \;=\; \mathrm{span}\{\ket{e_1},\dots,\ket{e_n}\}
\quad\text{(one excitation per block)},
\]
together with Hamiltonians $H_{\mathrm{pen}}$ and $H_{\mathrm{obj}}$ on $\OH$
such that:
\begin{enumerate}[label=\textup{(\alph*)}, leftmargin=2.2em]
\item \emph{Penalty structure.} $H_{\mathrm{pen}}$ is a sum of squared affine
      one–hot/degree/capacity penalties (optionally plus linear forbids) with
      integer coefficients bounded by $\mathrm{poly}(n)$. Consequently,
      $\mathrm{spec}(H_{\mathrm{pen}})\subseteq\{0,1,\dots,t_{\max}\}$ with
      $t_{\max}=O(m)=\mathrm{poly}(n)$.
\item \emph{Pattern symmetry.} $H_{\mathrm{pen}}$ is invariant under
      (i) block permutations $S_m$ and (ii) global symbol relabelings $S_n$.
      Hence the configuration space decomposes into level sets
      $L_t=\{x:\, H_{\mathrm{pen}}(x) :=\langle x \mid H_{\mathrm{pen}} \mid x \rangle \;=t\}$ that are preserved setwise.
\item \emph{Mixer and initial state.} The mixer family is block–local and normalized XY Hamiltonian,
      \[
      \widetilde H_{XY}^{(b)} \;=\;
      \frac{1}{n-1}\sum_{1\le a<b\le n}(X_aX_b+Y_aY_b),
      \qquad
      U_M(\beta) \;=\; \bigotimes_{b=1}^{m} e^{-i\beta\,\widetilde H_{XY}^{(b)}},
      \]
      with $\|\widetilde H_{XY}^{(b)}\|=O(1)$ on each block. The initial state is the
      uniform one–hot product
      \[
      \ket{s_0}\;=\;\ket{s_{\mathrm{blk}}}^{\otimes m},
      \qquad
      \ket{s_{\mathrm{blk}}}\;=\;\frac{1}{\sqrt n}\sum_{j=1}^n \ket{e_j}
      \quad\text{(a $W_n$ state initialized per block)}.
      \]
\end{enumerate}

\end{definition}

Evidently, the identifiable problem classes is broad and far from trivial. They include: Travelling Salesman (TSP/ATSP)~\cite{Lawler1985TSP,GareyJohnson1979}, Quadratic Assignment Problem (QAP)~\cite{Loiola2007QAPSurvey}, the shared transportation problems\cite{onah2025waas}, Generalized Assignment / Multiple–Knapsack~\cite{SahniGonzalez1976GAP}, $k$D Matching (NP-complete for $k\ge3$)~\cite{Karp1972}, etc. Some base cases are polynomial-time solvable but our interest is in the NP-hard variants (e.g. TSP/GAP/CVRP, etc.) for which the same one-hot encoding and symmetries apply.






\section{Results}
\label{sec:main}

\subsection{Fundamental Limitations}
\label{sec:upper}







Our point of departure is the following question. By how much can a shallow variational quantum circuit concentrate probability mass on the feasible manifold of a globally constrained problem?  We answer this for the canonical case of a permutation constrained problem of size $n$ one-hot encoded into $N=n^2$ qubits  (e.g., assignment/TSP) by proving that \emph{generic} QAOA initialized on \(|+\rangle^{\otimes N}\) with a $X$ mixer and a diagonal cost cannot raise the feasible mass above the uniform baseline \(|\Pi|/2^{N}\) by more than polynomial factors at any fixed constant depth, and under standard locality, not even at sublinear depth \(p=\alpha n\). Where $|\Pi|=n!$ and $\alpha < 1.$





\medskip
\noindent
Recall that the problem Hamiltonian is diagonal in computational basis and is denoted as $H_C$. For any $r$–local diagonal $H_C$ and all angles $(\beta,\gamma)$, we show that single layer $p=1$ generic QAOA circuits cannot generate the long–range correlations required by the feasible solution space by more than polynomial slack:
\begin{equation}
\label{eq:probs}
  P_{\Pi}^{\mathrm{(gen)}}(\beta,\gamma)\ \le\ \frac{|\Pi|}{2^{N}}\cdot \mathrm{poly}(n).
\end{equation}
The reason is twofold. Structurally, the $p{=}1$ light cone confines each block one–hot projector to a finite neighborhood, so the $O(n)$ constraints interact only through short dependencies. This caps any joint boost to at most polynomial over uniform. Spectrally, the indicator $1_\Pi$ is overwhelmingly high-degree in Walsh/Krawtchouk space, while a single $X$-mixer layer only adds phases. Consequently, average over $\gamma$ pins the baseline exactly and fourth-moment bounds show typical angles concentrate near it.


In contrast, a \emph{constraint-enhanced} variant (CE-QAOA) that operates inside a block one-hot subspace and mixes with an all-to-all \(XY\) block mixer achieves an \(\exp[\Theta(n^2)]\)  improvement over generic QAOA at the same depth. We arrive at the conclusion by establishing an upper envelope for generic QAOA and a lower envelope for CE--QAOA. To derive the upper envelopes, we employ harmonic analysis, angle-averaging, hyperconstrative fourth-moment, and light-cone analyses. Detailed proofs are presented in Methods and Materials ( Section \ref{sec:methods}).  We now summarize the main conclusions. %The different bounds answer different complementary questions.



\paragraph{Setup}
We refer to a COP of size $n$ as permutation constrained if the valid solutions can be identified as a $n \times n$ permutation matrix. For such problems, the one-hot encoding requires $N=n^2$ qubits. Let $N$ be the number of qubits and $\Pi=n!$ the size of the feasible permutation sector (row/column one–hot constrained). Consider the Boolean cube $\{0,1\}^N$ and write any function $f:\{0,1\}^N\!\to\!\mathbb{C}$ in the Walsh/Krawtchouk basis
$f(x)=\sum_{S\subseteq[N]}\widehat f(S)\,\chi_S(x)$ with characters $\chi_S(x)=(-1)^{\langle S,x\rangle}$ and \emph{degree} $|S|$. This is a complete and unique Fourier expansion for any function $f(x)$ with unique coefficients $\widehat f(S)$\cite{ODonnell2014}.

\medskip
\noindent
To connect this formalism the probability amplitudes in generic QAOA, we write
\[
  H_C=\sum_{y\in\{0,1\}^N} C(y)\,|y\rangle\!\langle y|
  \quad\Rightarrow\quad
  U_C(\gamma)=e^{-i\gamma H_C}
  =\sum_{y} e^{-i\gamma C(y)}\,|y\rangle\!\langle y|.
\]
With the phase field $\phi_\gamma(y):=e^{-i\gamma C(y)}$  and any basis vector $|x\rangle$ we obtain the probability profile
\[
  a_{\beta,\gamma}(x)
  =\langle x|U_X(\beta)U_C(\gamma)|+\rangle^{\otimes N}
  =2^{-N/2}\sum_{y} \underbrace{\langle x|U_X(\beta)|y\rangle}_{K_\beta(x,y)}\,\phi_\gamma(y).
\]
Because $U_X(\beta)=\prod_{i=1}^N e^{-i\beta X_i}$, its $Z$-basis kernel factors as
\begin{equation}
  K_\beta(x,y)
  \;=\; \prod_{i=1}^N \langle x_i|e^{-i\beta X}|y_i\rangle
  \;=\; (\cos\beta)^{\,N-d(x,y)}(-i\sin\beta)^{\,d(x,y)}
  \;=\; (\cos\beta)^{N}\,(-i\tan\beta)^{\,d(x,y)},
\end{equation}

where $d(x,y)=|\;x\oplus y\;|$ is the Hamming distance.  

\medskip
\noindent
Thus, we obtain a generic representation of the probability of obtaining a feasible permutation from generic QAOA as
\begin{equation}
\label{eq:generic-probability}
  a_{\beta,\gamma}(x)
  \;=\; 2^{-N/2}\sum_{y\in\{0,1\}^N} K_\beta(x,y)\,\phi_\gamma(y).
\end{equation}

For the indicator $1_\Pi$ on the permutation submanifold, we show (Materials and Methods Sections \ref{sec:walsh} and \ref{sec:kraw} ) that the total probability mass is tiny up to logarithmic degree. To derive this we define the \emph{low-degree mass} up to level $d$ as
\[
\|f\|^2_{\le d}\;:=\;\sum_{|S|\le d}\!|\widehat f(S)|^2,
\qquad
\text{and high-degree mass } \|f\|^2_{> d}= \|f\|_2^2-\|f\|^2_{\le d}.
\]
%  In our setting, $K_\beta$ \emph{diagonalizes} in the Walsh basis with shell phase
% $\lambda_\beta(|S|)=e^{-i\beta(N-2|S|)}$ (Lemma~\ref{lem:walsh-diag}), hence
% \begin{equation}
% \widehat{(K_\beta f)}(S)=\lambda_\beta(|S|)\,\widehat f(S),\qquad |\lambda_\beta(|S|)|=1.
% \end{equation}
% We call a layer \emph{degree–attenuating} if $|\lambda(|S|)|$ \emph{shrinks} as $|S|$ grows\cite{ODonnell2014,Beckner1975}. In our case, at $p{=}1$ there is \emph{no} attenuation because $|\lambda_\beta(|S|)|\equiv 1$ (pure phases). 
% Thus the $X$–mixer cannot reduce high–degree magnitudes; it only rephases shells.

\medskip
\noindent
We introduce Krawtchouk analysis to calculate the low–degree weight in the one-hot sector (Lemma~\ref{lem:onehot-lowdeg}), and, after tensorizing rows/columns and convolving to obtain the weight of indicator $1_\Pi$ on the permutation submanifold (Section \ref{sec:route1}). Our calculations show that both low–degree and high-degree mass remain exponentially suppressed. We denote the total output pmf as $p(x)=|a_{\beta,\gamma}(x)|^2$.  To compute the total feasible mass, we choose a degree truncation  $d=T=\Theta(\log n)$ and split $p(x)$ as low and high degree parts follows $p=p^{\le 2T}+p^{>2T}$. The low degree part accounts for all the contributions to the Fourier mass from degrees below the $T$ threshold and the high degree split captures all degree contributions above $T$. For the \emph{low–degree} part we have (Lemma~\ref{lem:lowdeg-corr-correct}),
\[
\frac{1}{2^N}\langle \mathbf{1}_\Pi, p^{\le 2T}\rangle
\le 
\frac{C_T n^{O(T)}}{2^{N}}.
\]




  
For the \emph{high–degree} tail is equally suppressed (Lemma~\ref{lem:tail-correct}). We have 
\[
\frac{1}{2^N}\big|\langle \mathbf{1}_\Pi, p^{>2T}\rangle\big|
\le \frac{\sqrt{|\Pi|}}{2^N}.
\]
Consequently, $p=p^{\le 2T}+p^{>2T}$ yields the harmonic bound in Theorem~\ref{thm:harmonic-baseline-correct}:
\[
P_\Pi(\beta,\gamma)\;\le\;\frac{|\Pi|}{2^N}.
\]
% \emph{Mechanism:} the only place where decay could enter is a degree–dependent $|\lambda(|S|)|<1$, but at $p{=}1$ the shell multipliers are unit-modulus phases, so the high–degree mass of $\mathbf{1}_\Pi$ remains undamped.

\medskip
\noindent
The angle averaging version of the result follows from an averaging over the phase field $\phi_\gamma(y)$ if the spectrum of $H_C$ is rescaled to lie on a lattice (after standard rescaling). This  arises because then $\mathbb{E}_\gamma[\phi_\gamma(y)\overline{\phi_\gamma(y')}]=\mathbf{1}_{\{C(y)=C(y')\}}$ leading cross–terms to cancel and
\[
\mathbb{E}_\gamma\big[P_\Pi(\beta,\gamma)\big]=\frac{|\Pi|}{2^N}.
\]
% Markov’s inequality gives high–probability concentration about the baseline for typical $\gamma$.

\medskip
\noindent
The third bound is derived the fourth–moment (hyperconractivity) analysis of the probability profile.  As outlined in Sec. \ref{app:L4}, let $a(x)=a_{\beta,\gamma}(x)$ and recall $P_\Pi=\sum_{x\in\Pi}|a(x)|^2$. 
By Cauchy–Schwarz,
\begin{equation}
  \sum_{x\in\Pi} |a_x|^2 \;\le\; \sqrt{|\Pi|}\,\Bigl(\sum_{x}|a_x|^4\Bigr)^{1/2}.
  \label{eq:CS-L4m}
\end{equation}
Thus, any bound on $\sum_x |a_x|^4$ yields an upper bound on $P_{\Pi}$.


Averaging over $\gamma$ (phase field), the one-bit calculation tensorizes to
\[
\mathbb{E}_\gamma\!\Big[\sum_x |a(x)|^4\Big]
\;\le\;\Big(\tfrac12+\tfrac14\sin^2(2\beta)\Big)^{\!N}
\;\le\;\Big(\tfrac34\Big)^{\!N},
\]
since $\max_{\beta}\big(\tfrac12+\tfrac14\sin^2 2\beta\big)=\tfrac34$ (achieved at $\beta=\pi/4$). 
Thus, for typical angles we have,
\[
P_\Pi(\beta,\gamma)
\;\lesssim\; \sqrt{|\Pi|}\,\Big(\tfrac34\Big)^{N/2}
\;=\; \frac{|\Pi|}{2^N}\cdot \underbrace{2^N\,|\Pi|^{-1/2}\,\Big(\tfrac34\Big)^{N/2}}_{\text{exponentially small in }N},
\]
i.e.\ exponentially close to the uniform baseline.

\medskip
\noindent
The light-cone analysis provides some structural intuition for the  above analysis. With $U=U_X(\beta)U_C(\gamma)$ and $H_C=\sum_j H_j$ diagonal, $r$–local, conjugation preserves diagonality and expands support by at most one $r$–local hyperedge. Each block projector $P_B$ depends only on its radius-1 neighborhood in the interaction hypergraph; the $2n$ one-hot constraints therefore form a bounded-dependency family. Standard cluster/dependency bounds then cap any feasible-mass boost to a polynomial over the uniform product baseline—fully consistent with the harmonic and moment bounds above.

\medskip
\noindent\textbf{Synthesis.}
All four routes—harmonic truncation with explicit low–degree suppression (Lemma ~\ref{lem:Pi-lowdeg} + Lemmas~\ref{lem:lowdeg-corr-correct},\ref{lem:tail-correct}), angle averaging via the phase field $\phi_\gamma$, fourth-moment with CS and the $(3/4)^N$ cap, and the light-cone locality barrier agree with Eq. \ref{eq:probs} with \emph{typical} angles exponentially close to $|\Pi|/2^N$. This is precisely the $p{=}1$ envelope used to compare against CE–QAOA, which evades the large $\frac{1}{2^N}$ suppression at any depth.




% Macros

\newcommand{\Perm}{\mathcal X}
\newcommand{\Pperm}{P_{\Perm}}
\newcommand{\UA}{\mathcal A}
% ===================== Setup =====================
\newcommand{\Hone}{\mathcal H_1}
\newcommand{\Htot}{\mathcal H}

\newcommand{\E}{\mathbb{E}}
% \newcommand{\ket}[1]{\left|#1\right\rangle}
% \newcommand{\bra}[1]{\left\langle#1\right|}
% \newcommand{\braket}[2]{\left\langle #1\,\middle|\,#2 \right\rangle}
% \newcommand{\ketbra}[2]{\left|#1\right\rangle\!\left\langle#2\right|}


\subsection{The Path to Exponential Enhancements}
%==================================================
% Depth-1 grid search on wi4 / wi5 (Aer, noiseless)
%==================================================
\begin{figure*}[t]
  \centering
  %--------------------------- wi4 ---------------------------
  \begin{subfigure}[t]{0.4\linewidth}
    \centering
    \includegraphics[width=\linewidth]{wi4_feasible_counts.png}
    \caption{\textbf{wi4, $n{=}4$ cities (16 qubits).}
             Blue bars: Block-QAOA ($p{=}1$) with \((\gamma,\beta)\) taken
             from the $10\times10$ grid.
             Orange bars (invisible) would correspond to generic QAOA,
             but every feasible bit-string occurred \emph{zero} times.}
    \label{fig:wi4_counts}
  \end{subfigure}
  %--------------------------- wi5 ---------------------------
  \begin{subfigure}[t]{0.4\linewidth}
    \centering
    \includegraphics[width=\linewidth]{wi5_feasible_counts.png}
    \caption{\textbf{wi5, $n{=}5$ cities (25 qubits).}
             Same colour code as (a).
             Again, generic QAOA produced no feasible outputs,
             while the constrained ansatz covers the full set of
             \(5!\!=\!120\) tours; the best tour appeared \(24\) times.}
    \label{fig:wi5_counts}
  \end{subfigure}
  \vspace{-0.7em}
  \caption{\textbf{Constraint-Enhanced QAOA dominates generic QAOA at
           depth \(p{=}1\).}   For medium-size TSP instances from QOptLib\cite{Osaba2024Qoptlib}, the unconstrained ansatz       spreads its amplitude almost entirely outside the            permutation subspace (\(n!/2^{n^{2}}\!\to\!0\)),           whereas the block-encoded circuit concentrates an
           \(\Omega(n^{-m})\) fraction of its probability on valid tours.}
  \label{fig:wi4_wi5_grid}
\end{figure*}


\begin{figure}[t]
  \centering
  \includegraphics[width=.65\linewidth]{ratio_emp.png}
  \caption{Empirical versus theoretical feasibility mass after parameter optimization on noiseless simulations at depth p = 1 for the three QOPLib benchmark instances wi3, wi4, and wi5 \cite{Osaba2024Qoptlib}. Each instance was executed with 500,000 shots so that Generic QAOA returns at least one feasible bit-string}.
    \label{fig:Blockqaoa_feasibility}
\end{figure}


 The \emph{constraint-enhanced} QAOA (CE-QAOA)  operates inside a product one-hot subspace and mixes with an all-to-all \(XY\) block mixer. Our main conclusion is that it achieves an \(\exp[\Theta(n^2)]\)  improvement over generic QAOA at the same depth. 
i.e.
\[\frac{P^{\mathrm{(CE)}}_{p}}{P^{\mathrm{(generic)}}_{p}}.\]
In order to make statements regarding the ratio $
\frac{P^{\mathrm{(CE)}}_{p}}{P^{\mathrm{(generic)}}_{p}}$, we need some robust lower envelope for CE--QAOA which we can compare to the upper envelope for generic QAOA outlined in Sec \ref{sec:upper}.  The necessary lower envelope comes from a basic symmetry of the product one-hot construction in Sec. \ref{sec:ce-qaoa}. For an instance in the kernel with $m$ blocks of local dimension $n$, a depth–$p$ CE–QAOA layer stack is given as 
\[
  \ket{\psi_p(\vec\gamma,\vec\beta)}
  \;=\;
  \Bigl(\prod_{\ell=1}^{p} U_M(\beta_\ell)\,e^{-i\gamma_\ell H_C}\Bigr)\ket{s_0},
  \qquad
  \vec\gamma=(\gamma_1,\dots,\gamma_p),\;
  \vec\beta=(\beta_1,\dots,\beta_p).
\]
Because $U_M$ preserves the one–hot sector and $H_C$ is diagonal, each layer maps the encoded manifold $\OH=(\mathcal H_1)^{\otimes m}$ to itself. The kernel’s pattern symmetry (block permutations $S_m$ and global symbol
relabelings $S_n$) permits a \emph{block–permutation twirl}. Averaging any fixed circuit $U$ over such relabelings yields a robust lower bound on the probability profile of the prepared state. For $p=1$, let \(\ket{x^\star}\in \OH\) be any fixed product basis vector and $D=n^n$ be the size of the product one-hot manifold explored by CE-QAOA, we recall the following lower bounds from Ref. \cite{onahce}:

 \begin{theorem}[Angle Agnostic Existence Bound \cite{onahce}] \label{thm:exist-params} For any product basis vector $\ket{x^\star}\in \OH$, and any angles \((\gamma,\beta)\), there exists a blockwise permutation \(\mathsf P^\star\) such that \[ \bigl|\langle x^{\star}\!\mid \mathsf P^{\star\dagger} U_{M}(\beta) U_{C}(\gamma) s_{0}\rangle \bigr|^{2} \;\ge\; \frac{c}{n^{m}}. \] \end{theorem} 


 \newcommand{\UM}{U_{M}}
\newcommand{\UC}{U_{C}}


 \begin{corollary}[Feasible–optimum specialization of Thm.~\ref{thm:exist-params}\cite{onahce}]
 \label{cor:feasible-optimum} Fix any grid angles $(\gamma,\beta)\in\mathcal G_{n+1}\times\mathcal G_{n+1}$ and set $U=\UM(\beta)\UC(\gamma)$ on $\OH$ with $D=\dim\OH=n^m$. Let $x^\star$ be any \emph{feasible optimal} label. Then there exists a blockwise permutation $\mathsf P^\star$ such that \[ \bigl|\langle x^\star \mid \mathsf P^{\star\dagger} U \mid s_0 \rangle\bigr|^2 \;\ge\; \frac{1}{D} \;=\; \frac{1}{n^m}. \] In particular, the constant in Theorem~\ref{thm:exist-params} can be taken as $c=1$. \end{corollary}

\medskip
\noindent




From the foregoing, we can therefore conclude that for every target basis
state \(x^{\star}\in\OH\), and in particular for every feasible optimal label,
the blockwise relabeling symmetry guarantees the existence of a blockwise
permutation \(\mathsf P^\star\) such that, for the fixed grid angles
\((\gamma,\boldsymbol\beta)\) under consideration,
\begin{equation}
  \bigl|\langle x^{\star}\!\mid
          \mathsf P^{\star\dagger}
          U_{M}(\boldsymbol\beta)
          U_{C}(\gamma)
          \psi_{0}\rangle
  \bigr|^{2}
  \;\ge\;
  \frac{1}{n^{m}}.
  \label{eq:lb}
\end{equation}
For permutation-constrained problems with \(m=n\), this gives the lower
envelope
\[
  P_{1}^{\mathrm{(CE)}} \;\ge\; n^{-n}.
\]

This lower bound is enough to prove exponential enhancements advertised in Eq \ref{eq:expp}. We organise the results in the following Theorem.




\begin{theorem}[Single Layer Exponential Enhancement] 
\label{thm:p1-sep}

Following the mixer and pattern symmetry in Def.~1, the blockwise relabeling
bound in Eq.~\eqref{eq:lb} gives the lower envelope
\[
  P_{1}^{\mathrm{(CE)}}\ge n^{-n}.
\]
For the generic QAOA, the feasible mass satisfies
\[
  \sup_{\beta,\gamma} P_{1}^{\mathrm{(generic)}}(\beta,\gamma)
  \le
  \frac{n!}{2^{n^2}}\cdot n^{O(1)}.
\]
Consequently,
\[
  \frac{P_{1}^{\mathrm{(CE)}}}{P_{1}^{\mathrm{(generic)}}}
  \ge
  \exp\!\left(n^2\ln 2-O(n\ln n)\right)
  =
  \exp[\Theta(n^2)].
\]
\end{theorem}

\begin{proof}
By Eq.~\eqref{eq:lb}, for permutation-constrained problems with \(m=n\), the
constraint-enhanced construction has the lower envelope
\[
  P_{1}^{\mathrm{(CE)}}\ge n^{-n}=\exp[-n\ln n].
\]
The generic QAOA upper envelope gives
\[
  \sup_{\beta,\gamma}P_{1}^{\mathrm{(generic)}}(\beta,\gamma)
  \le
  \frac{n!}{2^{n^2}}\cdot n^{O(1)}.
\]
Using \(\ln(n!)=n\ln n-O(n)\), we obtain
\[
  \log P_{1}^{\mathrm{(generic)}}
  \le
  -n^2\ln 2+n\ln n+O(n\ln n).
\]
Therefore
\[
  \log\frac{P_{1}^{\mathrm{(CE)}}}{P_{1}^{\mathrm{(generic)}}}
  \ge
  n^2\ln 2-2n\ln n-O(n\ln n)
  =
  \Theta(n^2).
\]
\end{proof}




\noindent
This supports the heuristic that to go beyond the near–uniform baseline on permutation–constrained problems one must \emph{build the constraints into the ansatz} (e.g., mixers that act within the target manifold). This is precisely the role of the mixer and pattern symmetry in the CE–QAOA Kernel which uses the problem's constraint structure as a resource and consequently yields a feasible–mass that sharply separates it from the generic ansatz.  



\subsection{Exponential Separation Beyond the Single layer}


The light-cone analysis can be extended to QAOA circuits beyond the single
layer in overlap-controlled sublinear-depth regimes.  In the one-dimensional
nearest-row geometry, the clean asymptotic condition is
\(p=o(n/\ln n)\), with constant-fraction \(n/\ln n\) windows depending on
the prefactors in the corrected exponent. For commuting, \(r\)-local diagonal costs, the Heisenberg light cone at depth \(p\) is confined to the radius-\(p\) neighborhood in the interaction hypergraph.  Projectors enforcing one-hot per row therefore depend on only \(O(p)\) (or \(O(\Delta^{p})\)) degrees of freedom. Where $\Delta^{p}$ upper-bounds the size of the radius-$p$ neighborhood of a vertex in a degree-$\Delta$ interaction hypergraph induced by the light-cone growth after $p$ layers. Using a special case of Finner/Hölder inequality on dependency graphs\cite{Finner1992} developed in Sec. \ref{app:fixed_p}, we obtain the following theorem.


\begin{theorem}[Exponential separation at constant/sublinear depth]
\label{thm:generic-const}
Let $N=n^2$ and $p=\alpha_n n$, where $\alpha_n\in(0,1)$ may depend on $n$.
Consider $p$–layers of the generic QAOA with X mixer and diagonal $H_C$.
Let $E_i$ denote the event that row $i$ is one--hot, and let $S_i(p)$ denote
the depth-$p$ row light cone of $E_i$. Define the row light-cone overlap number
\[
  \chi_p
  :=
  \max_{v\in[N]}
  \left|\{\,i\in[n]: v\in S_i(p)\,\}\right|.
\]
\begin{enumerate}
\item[(A)] \textbf{1D rows, nearest–row coupling.}
In this case \(W_{\rm row}(p)=O(2p+1)\), and
\[
  P_{p}^{\mathrm{(gen)}}
  \;\le\;
  \Bigl[C(2p{+}1)\,2^{-(n-1)}\Bigr]^{n/\chi_p}.
\]
In particular, since \(\chi_p\le 2p+1\),
\[
  P_{p}^{\mathrm{(gen)}}
  \;\le\;
  \Bigl[C(2p{+}1)\,2^{-(n-1)}\Bigr]^{n/(2p+1)}.
\]

\item[(B)] \textbf{Arbitrary intra–row degree $\Delta_{\rm row}$.}
In general,
\begin{equation}
  P_{p}^{\mathrm{(gen)}}
  \;\le\;
  \Bigl[C\,\Delta_{\rm row}^{\,p}\,2^{-(n-1)}\Bigr]^{n/\chi_p},
  \qquad
  \chi_p\lesssim \min\{n,\Delta_{\rm row}^{\,p}\}.
  \label{eq:linearn}
\end{equation}
Equivalently,
\[
  \log P_{p}^{\mathrm{(gen)}}
  \;\le\;
  -\frac{n}{\chi_p}
  \left[
    (n-1)\ln 2
    -
    \ln C
    -
    p\ln \Delta_{\rm row}
  \right].
\]
\end{enumerate}
Combining either (A) or (B) with the CE–QAOA lower bound
$P_{1}^{\mathrm{(CE)}}=\Omega(n^{-n})$ yields an exponential separation
whenever
\[
  \frac{n}{\chi_p}
  \left[
    (n-1)\ln 2
    -
    \ln\!\bigl(CW_{\rm row}(p)\bigr)
  \right]
  -
  n\ln n
  \;\longrightarrow\;+\infty .
\]
For the 1D nearest–row case, this includes the sublinear window
\[
  p=o\!\left(\frac{n}{\ln n}\right),
  \qquad\text{equivalently}\qquad
  \alpha_n=o\!\left(\frac{1}{\ln n}\right).
\]
Namely, a single layer of CE-QAOA can outperform $p$ layers generic QAOA
throughout the overlap-controlled sublinear-depth regime.

\end{theorem}

To prove this Theorem we shall need the following Lemma.



\begin{lemma}[Counting loss from row feasibility to permutation feasibility]
\label{lem:row-to-permutation-counting}
Let
\[
  \mathcal R
  :=
  \{x\in\{0,1\}^{n^2}: \text{each row of }x\text{ is one-hot}\}
\]
and let \(\Pi\subset\mathcal R\) be the set of permutation matrices, i.e.,
the set of bitstrings satisfying both row and column one-hot constraints.
Then
\[
  |\mathcal R|=n^n,
  \qquad
  |\Pi|=n!,
\]
and therefore
\[
  \frac{|\Pi|}{|\mathcal R|}
  =
  \frac{n!}{n^n}
  =
  \exp[-n+O(\log n)].
\]
Consequently, passing from row feasibility to full row-and-column
permutation feasibility contributes only an additional \(e^{-\Theta(n)}\)
counting factor under the uniform measure on \(\mathcal R\).
\end{lemma}

\begin{proof}
For each of the \(n\) rows, there are \(n\) possible positions for the unique
one, hence \(|\mathcal R|=n^n\).  Full row-and-column feasibility is exactly
the condition that the selected column labels form a permutation of
\([n]\), hence \(|\Pi|=n!\).  Stirling's formula gives
\[
  n!
  =
  \sqrt{2\pi n}\left(\frac ne\right)^n(1+o(1)),
\]
and therefore
\[
  \frac{n!}{n^n}
  =
  \sqrt{2\pi n}\,e^{-n}(1+o(1))
  =
  \exp[-n+O(\log n)].
\]
\end{proof}



\subsubsection{Proof of Theorem \ref{thm:generic-const}}% (Depth-\texorpdfstring{$p$}{p} success-probability bounds)}
\label{subsec:depth-p-bounds}

\paragraph{Setup.}Recall that the baseline uses the transverse–field mixer
$U_X(\beta)=\prod_i e^{-i\beta X_i}$ and a diagonal cost
$U_C(\gamma)=e^{-i\gamma H_C}$; thus light–cone growth is caused only by $U_C(\gamma)$.
Let $\Pi$ be the set of feasible bitstrings. For a fixed row (block) of length $n$, denote by $E_i$ the event “row $i$ is one–hot.” By the locality–parametrized light–cone lemma proved in the appendix
(Lemma~\ref{lem:row-valid-general}), in the 1D intra–row chain with nearest–row
coupling one has
\[
  \Pr[\text{row valid}]
  \;\le\;
  (2p{+}1)\,2^{-(n-1)}.
\]
More generally, for bounded intra–row degree $\Delta_{\rm row}$,
Lemma~\ref{lem:row-valid-general} gives
\(
  \Pr[\text{row valid}] \le C\,\Delta_{\rm row}^{\,p}\,2^{-(n-1)}
\)
for a constant $C$ depending only on the degree bounds. 

\medskip
\noindent 
Our light–cone lemmas are defined for the row–induced interaction graph, counting only count intra–row couplings which is usually slightly less than the full maximum degree. Informally, the Heisenberg support of a row one–hot projector after depth $p$
is confined to the radius–$p$ neighborhood in the interaction hypergraph; the
number of admissible completions inside this window scales as
$W_{\mathrm{row}}(p)\le 2p{+}1$ in 1D and $W_{\mathrm{row}}(p)=O(\Delta_{\rm row}^{\,p})$ in general, yielding the stated per–row probabilities which we now use in our proof.%~\cite{LiebRobinson1972,Hastings2010Locality}.

\begin{proof}[Sketch of Proof of Theorem \ref{thm:generic-const}]
For (A) and (B), insert the row–wise bounds from Lemma~\ref{lem:row-valid-general}
into the row events $E_i=\{\text{row $i$ is one–hot}\}$ and apply the chaining/
dependency-graph argument detailed in Sec.~\ref{app:fixed_p},
Thm.~\ref{thm:generic-const-app}. The relevant overlap number is
\[
  \chi_p
  :=
  \max_{v\in[N]}
  \left|\{\,i\in[n]: v\in S_i(p)\,\}\right|.
\]

Since full permutation feasibility implies row feasibility,
\[
  P_{\Pi,p}^{(\mathrm{gen})}
  =
  \Pr[X\in\Pi]
  \le
  \Pr\!\left[\bigcap_{i=1}^n E_i\right]
  =
  \mathbb E\!\left[\prod_{i=1}^n E_i\right].
\]
Applying the Finner/Hölder inequality~\eqref{eq:finner} gives
\[
  P_{\Pi,p}^{(\mathrm{gen})}
  \le
  \prod_{i=1}^n
  \bigl(\mathbb E[E_i^{\chi_p}]\bigr)^{1/\chi_p}.
\]
Since each \(E_i\) is an indicator, \(E_i^{\chi_p}=E_i\), and therefore
\[
  P_{\Pi,p}^{(\mathrm{gen})}
  \le
  \prod_{i=1}^n
  \bigl(\mathbb E[E_i]\bigr)^{1/\chi_p}.
\]

Substituting the per-row estimates yields the stated overlap-corrected bounds. The comparison with $P_{1}^{\mathrm{(CE)}}=\Omega(n^{-n})$ gives the stated
separation condition.
\end{proof}


\begin{remark}
Applying Finner~\eqref{eq:finner} directly to both rows and columns gives
\[
  P_{\Pi,p}^{(\mathrm{gen})}
  \le
  \prod_{i=1}^n
  \bigl(\mathbb E[R_i]\bigr)^{1/\widetilde\chi_p}
  \prod_{j=1}^n
  \bigl(\mathbb E[C_j]\bigr)^{1/\widetilde\chi_p}.
\]
If row and column validity obey the same light-cone bound, then
\[
  P_{\Pi,p}^{(\mathrm{gen})}
  \le
  \left[
    C W(p)2^{-(n-1)}
  \right]^{2n/\widetilde\chi_p}.
\]
The corresponding CE--vs--generic separation condition becomes
\[
  \frac{2n}{\widetilde\chi_p}
  \left[
    (n-1)\ln 2-\ln\!\bigl(CW(p)\bigr)
  \right]
  -
  n\ln n
  \to +\infty.
\]
\end{remark}



\subsection{Discussion}
Note that full permutation feasibility implies row feasibility.  Hence in the light-cone argument, one upper bounds \(P_\Pi^{(\mathrm{gen})}\) by the probability that all row one-hot events occur:
\[
  P_\Pi^{(\mathrm{gen})}
  \le
  \Pr\!\left[\bigcap_{i=1}^n E_i\right].
\]
The column constraints are not treated as independent additional events;
omitting them only weakens the upper bound.  The Finner/Hölder overlap
number \(\chi_p\) reduces the effective number of row constraints from \(n\)
to \(n/\chi_p\).

\begin{remark}[Depth monotonicity of CE--QAOA feasible mass]
\label{rem:ce-depth}
Throughout, let $P_p^{\mathrm{(CE)}}$ denote the \emph{angle-optimised}
feasible mass of depth-$p$ CE--QAOA, i.e.\ the maximum probability of
sampling a feasible string over all choices of $(\beta_1,\gamma_1,\dots,\beta_p,\gamma_p)$.
By construction, any depth-$(p{+}1)$ circuit can reproduce a depth-$p$ circuit by setting the final angles to zero, so the sequence is depth-monotone:
\[
  P_{p+1}^{\mathrm{(CE)}} \;\ge\; P_p^{\mathrm{(CE)}}\qquad\text{for all }p\ge 1.
\]
In particular, the one-layer lower bound
$P_{1}^{\mathrm{(CE)}} = \Omega(n^{-n})$ implied by Theorem~\ref{thm:exist-params} propagates to all deeper layers:
\[
  P_p^{\mathrm{(CE)}} \;\ge\; P_{1}^{\mathrm{(CE)}}
  \;=\; \Omega(n^{-n})
  \qquad\text{for every } p\ge 1.
\]

Combining this monotonicity with Theorem~\ref{thm:generic-const} yields a feasible-mass separation throughout any depth window in which the overlap-corrected generic exponent dominates the encoded-manifold baseline cost \(n\ln n\). In particular, for the one-dimensional nearest-row geometry this includes
\(p=o(n/\ln n)\).  Further sharpening of how \(P_p^{(\mathrm{CE})}\) itself grows with \(p\) is left for future work.
\end{remark}


\paragraph{Relation to known constant-depth limitations and OGP barriers.}
The above locality mechanism is consistent with several independent limitations:
(i) constant-depth circuits exhibit a bounded correlation radius/light-cone \cite{Hastings2010Locality} and hence cannot \emph{see} the whole instance structure at level~$p{=}1$ in the sense formalized for QAOA \cite{FarhiGamarnikGutmann2020WholeGraph,BravyiGossetKonig2018};
(ii) for random sparse graphs and related ensembles, the solution space exhibits the \emph{Overlap Gap Property} (OGP), which obstructs variational optimization at fixed depth \cite{Anschuetz_2024, GamarnikZadik2022QAOAOGP, FarhiGamarnikGutmann2020WholeGraph, Bravyi_2020}.
Formally, prior work shows (a) constant-level QAOA cannot surpass certain approximation thresholds on large sparse graphs \cite{Basso2022, Bravyi_2020}, and (b) under OGP/NLTS-type structure, low-depth circuits fail to reach near-optimal solutions with nonnegligible probability \cite{GamarnikZadik2022QAOAOGP,Anschuetz_2024}.



\paragraph{Row–level degree versus global degree.}
It will be important to distinguish the global hypergraph degree $\Delta$ of the
$N$-qubit interaction graph from the \emph{row–level} degree $\Delta_{\rm row}$.
By definition, $\Delta_{\rm row}$ is the maximum number of cost terms $H_a$ whose
support intersects a fixed row block
\[
  \mathrm{Row}_i \;=\; \{(i,1),\dots,(i,n)\},
\]
i.e., the maximum degree of the interaction graph induced on a single block of
$n$ qubits. In particular, $\Delta_{\rm row}$ depends only on the pattern of
\emph{intra–row} couplings and is independent of the total number of rows. In
our light–cone analysis, the Heisenberg evolution of a row one–hot projector is
confined to a neighborhood whose size grows at most like $O(\Delta_{\rm row}^{\,p})$,
so the number of admissible completions and hence the per–row validity
probability are controlled solely by $\Delta_{\rm row}$ (as made precise in
Lemma~\ref{lem:row-valid-general}), rather than by any global degree parameter
of the full $N$-qubit constraint graph. This row–level degree is therefore the
relevant locality parameter that ultimately bounds the joint feasibility
probability.








\paragraph{Non--asymptotic degree dependence in Eq.~\eqref{eq:linearn}.}
Let \(\Delta_{\rm row}(n)\) denote the maximum intra--row degree of the
interaction hypergraph induced by
\[
  H_C=\sum_a H_a(Z),
\]
where each \(H_a\) is diagonal, commuting, and \(r\)-local.  The row-validity
lemma gives
\[
  \mathbb E[E_i]
  \le
  C\,\Delta_{\rm row}(n)^p\,2^{-(n-1)}.
\]
The dependency--graph Finner/Hölder bound with row light-cone overlap number
\[
  \chi_p
  :=
  \max_{v\in[N]}
  \left|\{\,i\in[n]:v\in S_i(p)\,\}\right|
\]
therefore yields
\[
  P_p^{(\mathrm{gen})}
  \le
  \left[
    C\,\Delta_{\rm row}(n)^p\,2^{-(n-1)}
  \right]^{n/\chi_p}.
\]
Equivalently,
\[
  \log P_p^{(\mathrm{gen})}
  \le
  -\frac{n}{\chi_p}
  \left[
    (n-1)\ln 2
    -
    \ln C
    -
    p\ln \Delta_{\rm row}(n)
  \right].
\]
Combining this with the one-layer CE--QAOA lower envelope
\[
  P_1^{(\mathrm{CE})}
  =
  \Omega(n^{-n})
  =
  \exp[-O(n\ln n)]
\]
gives
\[
  \log
  \frac{P_1^{(\mathrm{CE})}}{P_p^{(\mathrm{gen})}}
  \ge
  \frac{n}{\chi_p}
  \left[
    (n-1)\ln 2
    -
    \ln C
    -
    p\ln \Delta_{\rm row}(n)
  \right]
  -
  O(n\ln n).
\]
Thus the CE--vs--generic ratio grows exponentially whenever the
overlap-corrected generic exponent dominates \(n\ln n\).  In the
one-dimensional nearest-row case, where
\[
  \chi_p\le 2p+1,
  \qquad
  W_{\rm row}(p)=O(2p+1),
\]
this includes
\[
  p=o\!\left(\frac{n}{\ln n}\right).
\]



\paragraph{Generic feasibility suppression.} The fully parametrized suppression is given as:



\[
  P_p^{(\mathrm{gen})}
  \;\lesssim\;
  \exp\!\left[
    -\frac{n}{\chi_p}
    \Bigl((n-1)\ln 2-\ln\!\bigl(CW_{\rm row}(p)\bigr)\Bigr)
  \right].
\]
Ignoring lower-order \(W_{\rm row}(p)\) terms, the suppression scale is
\[
  P_p^{(\mathrm{gen})}
  \;\lesssim\;
  \exp\!\left[
    -\Theta\!\left(\frac{n^2}{\chi_p}\right)
  \right].
\]
Consequently,
\[
  \chi_p=O(1)
  \quad\Longrightarrow\quad
  P_p^{(\mathrm{gen})}
  \le
  \exp[-\Theta(n^2)],
\]
\[
  \chi_p=o(n)
  \quad\Longrightarrow\quad
  P_p^{(\mathrm{gen})}
  \le
  \exp[-\omega(n)],
\]
and
\[
  \chi_p=\Theta(n)
  \quad\Longrightarrow\quad
  P_p^{(\mathrm{gen})}
  \le
  \exp[-\Theta(n)].
\]
Thus the generic feasibility suppression is governed by the effective
exponent \(n^2/\chi_p\).  The obstruction is quadratic-scale when
\(\chi_p=O(1)\), remains superlinear whenever \(\chi_p=o(n)\), and collapses
to only \(\exp[-\Theta(n)]\) once the light-cone overlap saturates at
\(\chi_p=\Theta(n)\).
For the CE--vs--generic comparison, the one-layer CE--QAOA lower envelope is
\[
  P_1^{(\mathrm{CE})}
  \ge
  n^{-n}
  =
  \exp[-n\ln n].
\]
Therefore the overlap-corrected generic exponent must dominate the
encoded-manifold baseline cost:
\[
  \frac{n^2}{\chi_p}
  \gg
  n\ln n.
\]






% ------------------------------------------------------------
%  Block- vs Generic-QAOA: success-probability gap (p = 1)


\section{Materials and Methods}
\label{sec:methods}




\subsection{Walsh–Fourier/Krawtchouk Analysis}
We seek to develop a Walsh–Fourier/Krawtchouk analysis that upper bounds the feasible mass of one–layer generic QAOA. Throughout, $N=n^2$, $\Pi\subset\{0,1\}^N$ is the feasible set of $n\times n$ permutation–matrix decoded from the sampled bitstrings, hence $|\Pi|=n!$ and 
\[
  \ket{\psi_1^{\mathrm{(gen)}}(\beta,\gamma)}
  \;=\; U_X(\beta)\,U_C(\gamma)\,\ket{+}^{\otimes N},
  \quad
  U_X(\beta) = e^{-i\beta\sum_{i=1}^N X_i},\quad
  U_C(\gamma)=e^{-i\gamma H_C}.
\]
With $H_C$ diagonal in the computational basis, we can write
\[
  H_C=\sum_{y\in\{0,1\}^N} C(y)\,|y\rangle\!\langle y|
  \quad\Rightarrow\quad
  U_C(\gamma)=e^{-i\gamma H_C}
  =\sum_{y} e^{-i\gamma C(y)}\,|y\rangle\!\langle y|.
\]
Acting on the uniform superposition gives
\[
  U_C(\gamma)|+\rangle^{\otimes N}
  =2^{-N/2}\sum_{y} e^{-i\gamma C(y)} |y\rangle.
\]
With the phase field $\phi_\gamma(y):=e^{-i\gamma C(y)}$  and any basis vector $|x\rangle$ we obtain
\[
  a_{\beta,\gamma}(x)
  =\langle x|U_X(\beta)U_C(\gamma)|+\rangle^{\otimes N}
  =2^{-N/2}\sum_{y} \underbrace{\langle x|U_X(\beta)|y\rangle}_{K_\beta(x,y)}\,\phi_\gamma(y).
\]
Because $U_X(\beta)=\prod_{i=1}^N e^{-i\beta X_i}$, its $Z$-basis kernel factors as
\begin{equation}
\label{eq:kb}
  K_\beta(x,y)
  \;=\; \prod_{i=1}^N \langle x_i|e^{-i\beta X}|y_i\rangle
  \;=\; (\cos\beta)^{\,N-d(x,y)}(-i\sin\beta)^{\,d(x,y)}
  \;=\; (\cos\beta)^{N}\,(-i\tan\beta)^{\,d(x,y)},
\end{equation}

where $d(x,y)=|\;x\oplus y\;|$ is the Hamming distance.  Thus the output amplitudes are
\begin{equation}
\label{eq:prob_def}
  a_{\beta,\gamma}(x)
  \;=\; 2^{-N/2}\sum_{y\in\{0,1\}^N} K_\beta(x,y)\,\phi_\gamma(y).
\end{equation}

For normalization, note that $K_\beta$ are matrix elements of the unitary $U_X(\beta)$, and compute
\begin{align}
  \sum_{x\in\{0,1\}^N} |a_{\beta,\gamma}(x)|^2
  &= 2^{-N}\sum_{x}\left|\sum_{y} K_\beta(x,y)\,\phi_\gamma(y)\right|^2
   \;=\; 2^{-N}\,\|\,U_X(\beta)\,\phi_\gamma\,\|_2^2 \\
  &= 2^{-N}\,\|\,\phi_\gamma\,\|_2^2
   \;=\; 2^{-N}\sum_{y} |\phi_\gamma(y)|^2
   \;=\; 2^{-N}\cdot 2^N \;=\; 1,
\end{align}
because $|\phi_\gamma(y)|\equiv 1$. Thus the amplitudes $a_{\beta,\gamma}(x)$ are correctly normalized.


\subsection{Walsh–Fourier primer}
\label{sec:walsh}


We write $[N]=\{1,2,\dots,N\}$. For a subset $S\subseteq[N]$, let $s\in\{0,1\}^N$
be its indicator vector, $s_i=\mathbf{1}_{\{i\in S\}}$, and we freely identify $S$ with $s$. The Walsh characters on the Boolean cube are~\cite{ODonnell2014}
\[
  \chi_S(x)\ :=\ (-1)^{S\cdot x},\qquad S\subseteq[N].
\]
For $x\in\{0,1\}^N$, define the mod-2 (parity) dot product
\[
  S\cdot x \;\coloneqq\; \Bigl(\sum_{i=1}^N s_i x_i\Bigr)\bmod 2 \;\in\; \{0,1\}.
\]

For $S,T\subseteq[N]$, write $S\oplus T$ for the symmetric difference
(bitwise XOR of indicators). Therefore:
\[
  S\oplus T \;\coloneqq\; (S\setminus T)\cup(T\setminus S),
\]
and 
\begin{align}
  (S\oplus T)\cdot x
  &\equiv (S\cdot x) + (T\cdot x) \pmod 2 \\
  &\equiv (S\cdot x)\oplus (T\cdot x).
\end{align}

So if $t$ is the indicator of $T$, then $S\oplus T$ corresponds to $s\oplus t$. The group law $\chi_S(x)\chi_T(x)=\chi_{S\oplus T}(x)$ will be used repeatedly. 



\medskip
\noindent
We use the (normalized) Walsh–Fourier transform is given as
\begin{align}
  \widehat{f}(S)\ :=\ 2^{-N}\!\sum_{x\in\{0,1\}^N} f(x)\,\chi_S(x),
  \qquad
  f(x)\ =\ \sum_{S\subseteq[N]} \widehat{f}(S)\,\chi_S(x).
  \label{eq:ft}
\end{align}
With this normalization, the hat $\widehat{\phantom{f}}$ denotes the Walsh–Fourier transform on the Boolean cube with characters
$\chi_S(x):=(-1)^{S\cdot x}$ for $S\subseteq[N]$.   Parseval/Plancherel reads~\cite{ODonnell2014}
\[
  \sum_{x}|f(x)|^2\ =\ 2^{N}\sum_{S}|\widehat{f}(S)|^2,
  \qquad
  2^{-N}\!\sum_{x}\overline{f(x)}\,g(x)\ =\ \sum_{S}\overline{\widehat{f}(S)}\,\widehat{g}(S).
\]




For functions $f:\{0,1\}^N\to\mathbb{C}$ we use the counting inner product
\(
  \langle f,g\rangle:=\sum_{x\in\{0,1\}^N}\overline{f(x)}\,g(x)
\)
(and the normalized version
\(
  \langle\!\langle f,g\rangle\!\rangle:=2^{-N}\langle f,g\rangle
\)
when convenient). 

\medskip
\noindent
One could view $K_\beta$ from Eq. \ref{eq:kb} as a convolution on $\mathbb{Z}_2^N$:
\[
  (K_\beta f)(x)\ :=\ \sum_{y} K_\beta(x,y)\,f(y)
  \ =\ \sum_{z} k_\beta(z)\,f(x\oplus z),
\]
with the one-bit kernel $k_\beta(0)=\cos\beta$ and $k_\beta(1)=-i\sin\beta$ and
$k_\beta(z)=\prod_{i=1}^N k_\beta(z_i)$.

\begin{lemma}[Diagonalization of $K_\beta$ in the Walsh basis]
\label{lem:walsh-diag}
For all $S\subseteq[N]$,
\[
  \widehat{(K_\beta f)}(S)\ =\ \lambda_\beta(|S|)\,\widehat{f}(S),
  \qquad
  \lambda_\beta(|S|)\ =\ e^{-i\beta\,(N-2|S|)}.
\]
In particular, $|\lambda_\beta(|S|)|=1$, so $K_\beta$ preserves Walsh magnitudes
$|\widehat{f}(S)|$.
\end{lemma}

\begin{proof}[Proof sketch]
On one bit, the convolution theorem with our normalization gives
$\widehat{k*f}(t)=2\,\widehat{k}(t)\,\widehat{f}(t)$ for $t\in\{0,1\}$, where
$\widehat{k}(t)=2^{-1}\sum_{z\in\{0,1\}}k(z)(-1)^{tz}$.
With $k(0)=\cos\beta$, $k(1)=-i\sin\beta$,
\[
  2\,\widehat{k}(0)=\cos\beta- i\sin\beta=e^{-i\beta},\qquad
  2\,\widehat{k}(1)=\cos\beta+ i\sin\beta=e^{i\beta}.
\]
By tensoring over $N$ bits, frequencies factorize. For a mask $S$ with $|S|$ ones,
$\lambda_\beta(|S|)=e^{-i\beta}$ on zero-frequencies and $e^{i\beta}$ on one-frequencies,
hence $\lambda_\beta(|S|)=e^{-i\beta(N-|S|)}e^{i\beta|S|}=e^{-i\beta(N-2|S|)}$.
\end{proof}

% \begin{tcolorbox}[title={Intuition: parity masks, phase shells, and feasible mass}, colback=white, breakable]
% \begin{itemize}
%   \item A subset $S\subseteq[N]$ is a \emph{parity mask}: $\chi_S(x)=+1$ if the number
%         of ones in $x$ on positions $S$ is even, and $-1$ if it is odd.
%   \item $K_\beta$ depends only on Hamming distance, so in Walsh space it acts
%         diagonally with phase $\lambda_\beta(|S|)=e^{-i\beta(N-2|S|)}$ per
%         “phase shell’’ $|S|$.
%   \item Since $|\lambda_\beta(|S|)|=1$, the $X$ rotation preserves $|\widehat f(S)|$;
%         feasible mass is driven by how $\mathbf{1}_\Pi$ correlates with the
%         Walsh autocorrelations of $U_C(\gamma)$’s phase field.
% \end{itemize}
% \end{tcolorbox}
\paragraph{Permutation–feasibility indicator.}
We write $\mathbf{1}_A:\{0,1\}^N\to\{0,1\}$ for the \emph{indicator} of a set $A$, i.e.
\[
  \mathbf{1}_A(x)=
  \begin{cases}
    1,& x\in A,\\
    0,& x\notin A.
  \end{cases}
\]
In the permutation encoding (e.g. for TSP) we fix $N=n^2$ and identify each bitstring
$x\in\{0,1\}^N$ with an $n\times n$ binary matrix
\[
  X(x)\;=\;\bigl(x_{i,j}\bigr)_{1\le i,j\le n},
\]
where $x_{i,j}\in\{0,1\}$ records whether city $j$ is placed at tour position $i$.
We say that $x$ is \emph{permutation–feasible} if and only if $X(x)$ is a
permutation matrix, i.e.
\[
  \sum_{j=1}^n x_{i,j} \;=\; 1\quad\text{for all rows }i\in[n],
  \qquad
  \sum_{i=1}^n x_{i,j} \;=\; 1\quad\text{for all columns }j\in[n].
\]
Let $\Pi\subset\{0,1\}^N$ denote the set of all such bitstrings, so
$|\Pi|=n!$. The permutation–feasibility indicator is
\[
  \mathbf{1}_\Pi(x)
  \;=\;
  \begin{cases}
    1, & \text{if }x\in\Pi\text{ (i.e.\ $X(x)$ is a permutation matrix)},\\[2pt]
    0, & \text{otherwise.}
  \end{cases}
\]
Equivalently, if we write
\[
  R(x)\;:=\;\prod_{i=1}^n \mathbf{1}\Bigl[\sum_{j=1}^n x_{i,j}=1\Bigr],
  \qquad
  C(x)\;:=\;\prod_{j=1}^n \mathbf{1}\Bigl[\sum_{i=1}^n x_{i,j}=1\Bigr],
\]
for the row– and column–one–hot indicators, then
\[
  \mathbf{1}_\Pi(x) \;=\; R(x)\,C(x),
  \qquad
  \|\mathbf{1}_\Pi\|_2^2 \;=\; \sum_{x}\mathbf{1}_\Pi(x) \;=\; |\Pi| \;=\; n!.
\]
With this notation, the feasible probability mass can be written as the Walsh correlation
\[
  P_\Pi(\beta,\gamma)
  \;=\; 2^{-N}\sum_{x}\mathbf{1}_\Pi(x)\,|g(x)|^2.
\]

Then
\[
  P_\Pi(\beta,\gamma)
  \ =\ 2^{-N}\sum_{x}\mathbf{1}_\Pi(x)\,|g(x)|^2
  \ =\ \sum_{S\subseteq[N]} \widehat{\mathbf{1}_\Pi}(S)\,\widehat{|g|^2}(S).
\]
where $g=K_\beta\phi_\gamma$ and $\widehat{\mathbf{1}_\Pi}$ is the Walsh transform
of the permutation–feasibility indicator. Writing $b:=\widehat{\phi_\gamma}$ and using Lemma~\ref{lem:walsh-diag},
\[
  \widehat{|g|^2}(S)
  \ =\ \sum_{T\subseteq[N]}\lambda_\beta(|T|)\,\overline{\lambda_\beta(|T\oplus S|)}\;
       b(T)\,\overline{b(T\oplus S)},
  \qquad
  \sum_{T}|b(T)|^2\ =\ 1,
\]
so $P_\Pi$ is controlled by the Walsh autocorrelations of the phase field and the Fourier profile of the feasible indicator.

%From Eq. \ref{eq:prob_def} $a(x)=a_{\beta,\gamma}(x)=2^{-N/2}g(x)$.


% Let $g:=K_\beta\phi_\gamma$ and $a=2^{-N/2}g$. Then, using Parseval/Plancherel on the Boolean cube,
% \begin{align}
%   P_\Pi(\beta,\gamma)
%   &= 2^{-N}\,\sum_{x\in\{0,1\}^N} \mathbf{1}_\Pi(x)\,|g(x)|^2
%    \;=\; \sum_{S\subseteq[N]} \widehat{\mathbf{1}_\Pi}(S)\,\widehat{|g|^2}(S).
%   \label{eq:mass-corr}
% \end{align}

% \section{Limitations of the Generic QAOA}
 
\subsection{Krawtchouk Polynomials on the Hamming sphere}
\label{sec:kraw}
% \paragraph{(I) Low–degree mass of one–hot/permutation indicators.}
Let $g_n$ be the “exactly–one” (one–hot) indicator on \(n\) bits, and let
\(\Pi\) be the intersection of \(2n\) block one–hot constraints plus the
row–column bijection (permutation) constraint. The Walsh spectrum of $g_n$ is explicit via Krawtchouk polynomials~\cite{feinsilver2007,ODonnell2014}. In particular,it allows us to show via explicit computation that low degrees carry
exponentially small mass. To see this, we work on the Boolean cube $\{0,1\}^n$ with Walsh characters $\chi_S(x)=(-1)^{S\cdot x}$ for $S\subseteq[n]$
and the normalization
\[
  \widehat f(S)\ :=\ 2^{-n}\sum_{x\in\{0,1\}^n} f(x)\,\chi_S(x),
  \qquad
  f(x)\ =\ \sum_{S\subseteq[n]} \widehat f(S)\,\chi_S(x)
\]
as in Eq. \ref{eq:ft}.

\medskip
\noindent
Let $g_w:\{0,1\}^n\to\{0,1\}$ denote the Hamming–sphere indicator $g_w(x)=\mathbf{1}_{\{|x|=w\}}$.
In particular, the one–hot indicator is $g_1$ (support $\{x:|x|=1\}$). For $0\le w\le n$, the (unnormalized) binary Krawtchouk polynomials are
\begin{equation}
  K_w^{(n)}(r)\ :=\ \sum_{j=0}^{w}(-1)^j \binom{r}{j}\binom{n-r}{\,w-j\,},
  \qquad r=0,1,\dots,n.
  \label{eq:Kraw-def}
\end{equation}
They satisfy the following generating function and orthogonality identities respectively~\cite{feinsilver2007}
\begin{align}
  \sum_{w=0}^{n} K_w^{(n)}(r)\, z^w \ &=\ (1+z)^{\,n-r}(1-z)^{\,r}, \label{eq:Kraw-gen}\\
  \sum_{r=0}^{n} \binom{n}{r}\, K_w^{(n)}(r)\,K_{w'}^{(n)}(r) \ &=\ 2^{n}\binom{n}{w}\,\delta_{w,w'}.
  \label{eq:Kraw-orth}
\end{align}

\begin{lemma}[Krawtchouk transform of a Hamming sphere]
\label{lem:layer-sum-kraw}
For $S\subseteq[n]$ with $|S|=r$,
\[
  \sum_{x:\,|x|=w} \chi_S(x) \;=\; K_w^{(n)}(r).
\]
Consequently,
\begin{equation}
  \widehat{g_w}(S) \;=\; 2^{-n}\,K_w^{(n)}\!\bigl(|S|\bigr).
  \label{eq:gw-hat}
\end{equation}
\end{lemma}

\begin{proof}
Partition the coordinates into the $r$ positions where $S$ has 1s and the $n-r$ positions where $S$ has 0s.
To form an $x$ with $|x|=w$, choose $j$ ones among the $r$ ``active'' positions and $w-j$ ones among the remaining $n-r$ positions.
There are $\binom{r}{j}\binom{n-r}{w-j}$ such $x$, and each contributes a phase $(-1)^j$ because $\chi_S(x)=(-1)^{S\cdot x}$ counts parity on the $r$ active positions only.
Summing over $j$ gives \eqref{eq:Kraw-def}, hence the first claim.
Applying the Walsh transform definition to $g_w$,
\[
  \widehat{g_w}(S)=2^{-n}\sum_{x} g_w(x)\chi_S(x) 
  = 2^{-n}\sum_{|x|=w}\chi_S(x)=2^{-n}K_w^{(n)}(|S|).
\]
\end{proof}

\[
  K_{1}^{(n)}(r)\ =\ \binom{r}{0}\binom{n-r}{1}-\binom{r}{1}\binom{n-r}{0}
  \ =\ (n-r)-r\ =\ n-2r.
\]
Thus, for $g_1$ (the ``exactly one'' indicator),
\begin{equation}
  \widehat{g_1}(S)\ =\ 2^{-n}\bigl(n-2|S|\bigr),
  \qquad S\subseteq[n],
  \label{eq:onehot-transform}
\end{equation}
which matches the direct computation $\widehat{g_1}(S)=2^{-n}\sum_{i=1}^{n}(-1)^{\mathbf{1}_{\{i\in S\}}}=2^{-n}(n-2|S|)$.

Normalization/Parseval sanity check  using \eqref{eq:gw-hat} and \eqref{eq:Kraw-orth} gives,
\[
  \sum_{x}|g_w(x)|^2 = \binom{n}{w}
  \ =\ 2^{n}\sum_{S}|\widehat{g_w}(S)|^2
  \ =\ 2^{n}\sum_{r=0}^{n}\binom{n}{r}\,\bigl(2^{-n}K_w^{(n)}(r)\bigr)^2,
\]
which is exactly \eqref{eq:Kraw-orth} with $w'=w$. We also have
$\|g_1\|_2^2=\binom{n}{1}=n$.

\begin{lemma}[Low–degree Fourier weight for $g_1$]
\label{lem:onehot-lowdeg}
For any fixed degree cutoff $d\ge 0$, where $d\in\mathbb{N}$ denotes a \emph{Fourier degree cutoff}
\[
  \sum_{\,|S|\le d}\! \bigl|\widehat{g_1}(S)\bigr|^2
  \ =\ 2^{-2n}\sum_{r=0}^{d} \binom{n}{r}\,(n-2r)^2
  \ \le\ 2^{-2n}\,C_d\,n^{\,d+2},
\]
for some constant $C_d$ depending only on $d$ (e.g.\ $C_d=2(d+1)$ for all $n\ge 2$).
Equivalently, the degree–$\le d$ weight is at most $C_d\,n^{\,d+1}\,2^{-n}$ of the total $\|g_1\|_2^2=n$.
\end{lemma}

\begin{proof}
Insert \eqref{eq:onehot-transform} and group by $r=|S|$:
\[
  \sum_{|S|\le d}\! \bigl|\widehat{g_1}(S)\bigr|^2
  = 2^{-2n}\sum_{r=0}^{d} \binom{n}{r}\,(n-2r)^2.
\]
For $r\le d$ we have $\binom{n}{r}\le n^{r}$ and $(n-2r)^2\le n^{2}$, hence
\[
  \sum_{r=0}^{d} \binom{n}{r}\,(n-2r)^2
  \ \le\ n^{2}\sum_{r=0}^{d} n^{r}
  \ =\ n^{2}\,\frac{n^{d+1}-1}{n-1}
  \ \le\ 2(d+1)\,n^{\,d+1}\qquad(n\ge 2).
\]
Multiplying by $2^{-2n}$ yields the claimed bound with $C_d=2(d+1)$.
\end{proof}


The Krawtchouk identity $\widehat{g_w}(S)=2^{-n}K_w^{(n)}(|S|)$ shows that the Walsh spectrum
of any Hamming–sphere indicator is a radial (shell–dependent) polynomial of $|S|$.
For the one–hot layer ($w=1$), $\widehat{g_1}(S)$ is exactly linear in $|S|$,
and its low–degree weight decays exponentially in $n$ after summation over shells, as quantified in Lemma~\ref{lem:onehot-lowdeg}. Tensorizing across $m=2n$ blocks and incorporating the bijection constraint yields an additional sparse factor that does not increase low–degree mass beyond a polynomial.


\subsection{Route 1: Harmonic Analysis of the probability profile}
\label{sec:route1}

To go beyond $n-$bit rows and make the Fourier weight decay on the permutation feasible subspace explicit, we shall use the following technical lemmas. Lemma \ref{lem:ratio-nn-over-nfact} is an adaptation from  Alon and Spencer (Lemma~2, p.~64, \cite{AlonSpencer2008}) while Lemma \ref{lem:control-ineq} is a technical tool to control the resulting inequality.

\begin{lemma}
\label{lem:ratio-nn-over-nfact}
For every integer $n\ge 2$,
\[
\bigl((n-1)!\bigr)^{\,n(n+1)}\cdot \bigl((n+1)!\bigr)^{\,n(n-1)}
\;<\;
\bigl(n!\bigr)^{\,2(n^2-1)},
\]
equivalently,
\[
\left(\frac{n^n}{n!}\right)^{\!2}
\;>\;
\left(\frac{n+1}{n}\right)^{\,n(n-1)}.
\]
\end{lemma}

\begin{proof}
The two forms are equivalent by dividing both sides of the first inequality by
$\bigl((n-1)!\,(n+1)!\bigr)^{\,n(n-1)}$ and simplifying.

\emph{Base case $n=2$.} We have
\[
\left(\frac{2^2}{2!}\right)^{\!2}= \left(\frac{4}{2}\right)^2=4
\quad\text{and}\quad
\left(\frac{3}{2}\right)^{2}= \frac{9}{4},
\]
so the inequality holds.

\emph{Case $n\ge 3$.} By Arithmetic Mean–Geometric Mean inequality ( AM--GM),
\[
n!\;=\;\prod_{k=1}^n k \;\le\; \left(\frac{1+\cdots+n}{n}\right)^{\!n}
=\left(\frac{n+1}{2}\right)^{\!n}.
\]
Hence
\[
\left(\frac{n^n}{n!}\right)^{\!2}
\;\ge\;
\left(\frac{n^n}{\bigl(\tfrac{n+1}{2}\bigr)^n}\right)^{\!2}
=
\left(\frac{2n}{n+1}\right)^{\!2n}.
\]
Using $(1+\tfrac{1}{n})^n<e$ we get
\[
\frac{2n}{n+1}
=\frac{2}{1+\tfrac{1}{n}}
> 2e^{-1/n},
\qquad\Rightarrow\qquad
\left(\frac{2n}{n+1}\right)^{\!2n} > 4^n e^{-2}.
\]
For $n\ge 3$ we have $4^n>e^{\,n+1}$, hence $4^n e^{-2}>e^{\,n-1}$. Therefore
\[
\left(\frac{n^n}{n!}\right)^{\!2}
> e^{\,n-1}.
\]
On the other hand, $(1+\tfrac{1}{n})^n<e$ implies
\[
\left(\frac{n+1}{n}\right)^{\,n(n-1)} = \bigl((1+\tfrac{1}{n})^n\bigr)^{\,n-1} < e^{\,n-1}.
\]
Combining the last two displays yields
\[
\left(\frac{n^n}{n!}\right)^{\!2}
\;>\;
\left(\frac{n+1}{n}\right)^{\,n(n-1)}.
\]
% as claimed.
\end{proof}

\begin{lemma}
\label{lem:control-ineq}
Let $c_T>0$ and set $a:=\tfrac{c_T\ln 2}{2}$. Then for all
\[
n \;\ge\; N(c_T)\ :=\ \max\!\Bigl\{9,\ \Bigl\lceil \tfrac{4}{a^2}\Bigr\rceil\Bigr\}
\]
we have $\,n^n\,2^{-c_T n^2/2}<1$.
\end{lemma}

\begin{proof}
For $n\ge 9$ we have $\ln n<\sqrt{n}$, and for $n\ge 4/a^2$ we have
$\sqrt{n}\le a n$. Hence $\ln n<a n$, so
\[
n\ln n-a n^2 \;<\; 0,
\]
which implies $n^n 2^{-c_T n^2/2}=\exp(n\ln n-a n^2)<1$.
\end{proof}


\begin{lemma}[Low–degree mass of $\mathbf{1}_\Pi$]
\label{lem:Pi-lowdeg}
Let $N=n^2$ and $\Pi\subset\{0,1\}^{N}$ be the set of $n\times n$ permutation
matrices (row– and column–one–hot). For every fixed degree cutoff $d\in\mathbb{N}$ there exist
constants $c_d,C_d>0$ such that
\[
  \sum_{|S|\le d}\,|\widehat{\mathbf{1}_\Pi}(S)|^2
  \;\le\; C_d\,n^{O(d)}\,2^{-c n^2}\,|\Pi|^2R_n;
  \qquad \|\mathbf{1}_\Pi\|_2^2=|\Pi|=n! \ \textit{and } R_n \;>\; \Bigl(\tfrac{n+1}{n}\Bigr)^{n(n-1)}
\]
\end{lemma}

\begin{proof}
Write the indicator of \emph{row} one–hot as
\[
  R(x)\;=\;\prod_{i=1}^{n} g^{(i)}_1\!\bigl(x_{i,1},\dots,x_{i,n}\bigr),
\]
where each $g_1^{(i)}$ is the (one–bit–weight) indicator on the $i$-th row block
(size $n$). Similarly, the \emph{column} one–hot indicator factors as
\[
  C(x)\;=\;\prod_{j=1}^{n} g^{(j)}_1\!\bigl(x_{1,j},\dots,x_{n,j}\bigr).
\]
Then $\mathbf{1}_\Pi(x)=R(x)\,C(x)$ and, in Walsh–Fourier space,
\[
  \widehat{\mathbf{1}_\Pi}\;=\;\widehat{R\,C}\;=\;\widehat{R} * \widehat{C},
\]
where $*$ is convolution on the frequency group $(\{0,1\}^N,\oplus)$:
$(f*g)(S)=\sum_{T} f(T)\,g(S\oplus T)$.

\medskip
\noindent\textbf{Step 1: Tensorization across blocks (row side).}
By block disjointness of the rows and the factorization of characters,
the Walsh transform \emph{factorizes}:
for $S\subseteq[N]$ and $S_i:=S\cap\{\text{positions of row }i\}$,
\[
  \widehat{R}(S)\;=\;\prod_{i=1}^{n} \widehat{g_1}\bigl(S_i\bigr).
\]
By Lemma~\ref{lem:onehot-lowdeg} (the Krawtchouk transform with $w=1$),
$\widehat{g_1}(T)=2^{-n}\bigl(n-2|T|\bigr)$ for any $T\subseteq[n]$. Note that since a single \(n\)-bit one–hot block \(g_1\) has \(\widehat{g_1}(T)=2^{-n}(n-2|T|)\).
Tensoring across \(n\) disjoint rows gives \(\|\widehat{R}\|_2^2=2^{-n^2}n^n\), and
similarly for columns \(\|\widehat{C}\|_2^2=2^{-n^2}n^n\). 
Hence
\begin{equation}
  \widehat{R}(S)
  \;=\;2^{-n^2}\prod_{i=1}^{n}\bigl(n-2|S_i|\bigr),
  \qquad
  \|\widehat{R}\|_2^2 \;=\; 2^{-N}\,\|R\|_2^2 \;=\; 2^{-N}\,n^n.
  \label{eq:Rhat-factor}
\end{equation}

An analogous factorization holds for $C$, partitioning $S$ by columns:
$\widehat{C}(S)=2^{-n^2}\prod_{j=1}^{n}\bigl(n-2|S^{(j)}|\bigr)$ and
$\|\widehat{C}\|_2^2=2^{-N}n^n$.

\medskip
\noindent\textbf{Step 2: Convolution bound onto a low–degree window.}
Let $E_d:=\{S\subseteq[N]:|S|\le d\}$ denote the degree–$\le d$ frequency window.
For any $S$,
\(
(\widehat{R}*\widehat{C})(S)=\sum_{T}\widehat{R}(T)\widehat{C}(S\oplus T).
\)
By Cauchy–Schwarz in the $T$–sum and then summing over $S\in E_d$,
\begin{align*}
  \sum_{S\in E_d} |(\widehat{R}*\widehat{C})(S)|^2
  &\le \sum_{S\in E_d}
        \Bigl(\sum_T |\widehat{R}(T)|^2\Bigr)
        \Bigl(\sum_T |\widehat{C}(S\oplus T)|^2\Bigr) \\
  &= \Bigl(\sum_T |\widehat{R}(T)|^2\Bigr)
     \Bigl(\sum_{S\in E_d}\sum_T |\widehat{C}(S\oplus T)|^2\Bigr) \\
  &= \|\widehat{R}\|_2^2\;\bigl|E_d\bigr|\;\|\widehat{C}\|_2^2
   \;=\; \bigl|E_d\bigr|\;\cdot\;2^{-2N}\,n^{2n},
\end{align*}
using $\sum_T |\widehat{C}(S\oplus T)|^2=\|\widehat{C}\|_2^2$ (shift invariance).
The window size satisfies, for fixed $d$, the standard bound
\[
  |E_d|=\sum_{d=0}^{d}\binom{N}{d}\;\le\; \Bigl(\frac{eN}{d}\Bigr)^{\!d}
  \;=\; N^{O(d)} \;=\; n^{O(d)}.
\]
Therefore
\begin{equation}
  \sum_{|S|\le d} |\widehat{\mathbf{1}_\Pi}(S)|^2
  \;\le\; C_d\,N^{O(d)}\;\cdot\;2^{-2N}\,n^{2n}
  \;=\; C_d\,n^{O(d)}\;2^{-2n^2}\,n^{2n}.
  \label{eq:lowdeg-Pi-pre}
\end{equation}

\medskip
\noindent\textbf{Step 3: using Lemma~\ref{lem:ratio-nn-over-nfact} on  \eqref{eq:lowdeg-Pi-pre}}
\[
\sum_{|S|\le d} |\widehat{\mathbf{1}_\Pi}(S)|^2
\;\le\; C_d\,n^{O(d)}\,2^{-2n^2}\,n^{2n}.
\]
Insert the exact factorization
\[
n^{2n} \;=\; (n!)^2\!\left(\frac{n^n}{n!}\right)^{\!2}
\;=\; (n!)^2\,R_n,
\qquad R_n:=\left(\frac{n^n}{n!}\right)^{\!2}.
\]
By Lemma~\ref{lem:ratio-nn-over-nfact} (equivalent form),
\[
R_n \;>\; \Bigl(\tfrac{n+1}{n}\Bigr)^{n(n-1)}
\;=\;\exp\!\bigl(\Theta(n)\bigr),
\]
so the multiplicative gap between $n^{2n}$ and $(n!)^2$ is only $\exp(\Theta(n))$.
Therefore,
\[
\sum_{|S|\le d} |\widehat{\mathbf{1}_\Pi}(S)|^2
\;\le\; C_d\,n^{O(d)}\,2^{-2n^2}\,(n!)^2R_n.
\]
Using $\|\mathbf{1}_\Pi\|_2^2=|\Pi|=n!$, we get% (for all sufficiently large $n$),
% \[
% \sum_{|S|\le d} |\widehat{\mathbf{1}_\Pi}(S)|^2
% \;\le\; C_d\,n^{O(d)}\,2^{-c n^2}\,|\Pi|^2R_n.
% \]

\[
\sum_{|S|\le d} |\widehat{\mathbf{1}_\Pi}(S)|^2
\;\le\; C_d\,n^{O(d)}\,2^{-c n^2}\,|\Pi|^2R_n.
\]

\end{proof}


The \(2^{-n}\) per block across \(n\) rows and \(n\) columns is the source of the \(2^{-\Theta(n^2)}\) decay. To bound the total probability profile, write the low/high–degree split of the contributuions to toal probability as \(p=p^{\le 2T}+p^{>2T}\), where \(p^{\le2T}\) denotes the Walsh truncation to degrees \(\le 2T\) and effectively identify $d=2T$. We use Lemma~\ref{lem:ratio-nn-over-nfact} to replace $n^{2n}$
by $(n!)^2 R_n$, exposing the $(n!)^2$ factor explicitly. $R_n$ remains suppressed by $2^{-c n^2}$. Additional suppression will come from the normalization terms as we shall see in the next Lemmas.

\begin{lemma}[Low–degree probability mass contribution]
\label{lem:lowdeg-corr-correct}
For every cutoff degree  $T\in\mathbb{N}$, there exist a constant $C_T>0$ s.t.
\[
  2^{-N}\,\bigl\langle \mathbf{1}_\Pi,\; p^{\le2T}\bigr\rangle
  \;\le\;
  \frac{C_T}{2^{N}}\; n^{O(T)} 
\]

\end{lemma}

\begin{proof}
By Plancherel,
\[
  2^{-N}\,\langle \mathbf{1}_\Pi, p^{\le 2T}\rangle
  = \sum_{|S|\le 2T}\widehat{\mathbf{1}_\Pi}(S)\,\overline{\widehat p(S)}
  \;\le\;
  \Bigl(\sum_{|S|\le2T}|\widehat{\mathbf{1}_\Pi}(S)|^2\Bigr)^{1/2}
  \Bigl(\sum_{|S|\le2T}|\widehat p(S)|^2\Bigr)^{1/2},
\]

 Using Lemma ~\ref{lem:Pi-lowdeg} in \emph{fractional-energy} form
\[
  \sum_{|S|\le2T}|\widehat{\mathbf{1}_\Pi}(S)|^2
  \le C_T n^{O(T)} 2^{-c_T n^2}\cdot|\Pi|^2R_n.
\]


For the second factor, use
$\sum_S|\widehat p(S)|^2 = 2^{-N}\|p\|_2^2 \le 2^{-N}$ since $p$ is a pmf.
Collecting terms gives


\[
  2^{-N}\,\bigl\langle \mathbf{1}_\Pi,\; p^{\le2T}\bigr\rangle
  \;\le\;
  \Bigl(C_T n^{O(T)} 2^{-c_T n^2}\tfrac{|\Pi|}{2^{N}}\Bigr)^{\!1/2}
  \Bigl(\tfrac{R_n\Pi}{2^{N}}\Bigr)^{\!1/2}
  \;\le\;
  \frac{{|\Pi|}}{2^{N}}\;  \sqrt{R_n}.C_T n^{O(T)} 2^{-c_T n^2/2}
\]


Recall that $ R_n:=\left(\frac{n^n}{n!}\right)^{\!2}$ so that $\sqrt{R_n} =\left(\frac{n^n}{n!}\right) = \frac{n^n}{|\Pi|}. $

Therefore:
\[
  2^{-N}\,\bigl\langle \mathbf{1}_\Pi,\; p^{\le2T}\bigr\rangle
  \;\le\;
  \frac{{n^n}}{2^{N}}\; .C_T n^{O(T)} 2^{-c_T n^2/2}
    \;\le\;
     \frac{{C_T}}{2^{N}}\; . n^{O(T)} 
\]
Where we have used the fact that \(n^n\cdot2^{-c_T n^2/2} <1\) from Lemma \ref{lem:control-ineq} in the last step.
\end{proof}

 \begin{lemma}[High–degree contribution]
\label{lem:tail-correct}
For all $T\in\mathbb{N}$ and all angles $(\beta,\gamma)$,
\[
  2^{-N}\,\bigl|\langle \mathbf{1}_\Pi,\; p^{>2T}\rangle\bigr|
  \;\le\; \frac{\|\mathbf{1}_\Pi\|_2}{2^N}\,\|p^{>2T}\|_2
  \;\le\; \frac{\sqrt{|\Pi|}}{2^{N}}.
\]
Moreover, without additional structure on $\phi_\gamma$, one cannot obtain a
uniform $T$–decaying bound here because the Walsh multipliers of $K_\beta$ satisfy
$|\kappa_\beta(t)|=1$ for all $t$ (pure phases), so there is no degree–attenuation at $p=1$.
\end{lemma}

\begin{proof}
Cauchy–Schwarz gives
$|\langle \mathbf{1}_\Pi, p^{>2T}\rangle|\le\|\mathbf{1}_\Pi\|_2\,\|p^{>2T}\|_2$.
Then normalize by $2^{N}$ to get the first inequality. Since $p^{>2T}$ is the orthogonal Walsh projection of $p$, $\|p^{>2T}\|_2\le \|p\|_2$.
Because $p$ is a pmf, $\|p\|_2^2=\sum_x p(x)^2\le \sum_x p(x)\|p\|_\infty=1$, hence $\|p^{>2T}\|_2\le 1$.
Also $\|\mathbf{1}_\Pi\|_2=\sqrt{|\Pi|}$.
For the remark, $\kappa_\beta(t)=e^{-i\beta(N-2t)}$ so $|\kappa_\beta(t)|=1$.
\end{proof}


Adding column one-hot (turning “row-one-hot” into permutation matrices) multiplies by a sparse indicator that enforces n extra constraints. In Fourier space that becomes a convolution that can at worst blow up the low-degree mass by a polynomial factor in n and d (Cauchy–Schwarz on the convolution over the degree simplex). 



\begin{theorem}[Harmonic–analysis baseline]
\label{thm:harmonic-baseline-correct}
There exist absolute constants $c,C>0$ such that taking the cutoff degree  $T\in\mathbb{N}$  with $T=\lceil c\log n\rceil$,
\[
  P_\Pi(\beta,\gamma)
  \;\le\;
  \frac{|\Pi|}{2^{N}}.
\]
\end{theorem}

\begin{proof}
Choosing $T=\lceil c\log n\rceil$ , split $p=p^{\le 2T}+p^{>2T}$. Then
\[
  P_\Pi = 2^{-N}\langle \mathbf{1}_\Pi, p^{\le 2T}\rangle
         + 2^{-N}\langle \mathbf{1}_\Pi, p^{> 2T}\rangle.
\]
\emph{Low–degree part.}
By Lemma~\ref{lem:lowdeg-corr-correct}, setting $C=C_T$,
\[
  2^{-N}\,\bigl\langle \mathbf{1}_\Pi,\; p^{\le2T}\bigr\rangle
    \;\le\;
     \frac{{C}}{2^{N}}\; . n^{O(T)} 
\]


\emph{High–degree part.}
By Lemma~\ref{lem:tail-correct},
\[
  2^{-N}\bigl|\langle \mathbf{1}_\Pi, p^{>2T}\rangle\bigr|
  \;\le\; \frac{\|\mathbf{1}_\Pi\|_2}{2^{N}}\;\|p^{>2T}\|_2
  \;\le\; \frac{\sqrt{|\Pi|}}{2^{N}}.
\]



Summing the two contributions gives
\[
  P_\Pi(\beta,\gamma)
  \;\le\;  \frac{{C}}{2^{N}}\; . n^{O(T)} 
           \;+\;  \frac{\sqrt{|\Pi|}}{2^{N}}
            \;=\; \frac{({C \cdot n^{O(T)} + \sqrt{|\Pi|})}}{2^{N}}\;
  \;\le\; \frac{|\Pi|}{2^{N}}.
\]
Where we have used that the numerator sums to a quantity less that $|\Pi|$ for any large $n$. This completes the proof.
\end{proof}




\subsection{Route II: Parameter–Averaging (Exact Expectation and a Typical–Angles Bound)}
\label{app:avg-angles}

In this section, we show that averaging over the \emph{cost angle} pins the baseline feasible mass to the uniform value $|\Pi|/2^{N}$, independently of the mixer angle, provided the cost spectrum lies on a lattice (which holds after rescaling for the usual integer/rational Ising/TSP costs). This insight comes simple Fourier-averaging. 

\medskip
\noindent
Recall that averaging over a \emph{wide} distribution of the cost angle $\gamma$ makes the oscillatory
phases $e^{-i\gamma\,\Delta}$ cancel out for every fixed nonzero frequency $\Delta$.
When we expand $\mathbb{E}_\gamma\!\left[|a_x(\beta,\gamma)|^2\right]$, the cross terms
carry factors $e^{-i\gamma\,[C(y)-C(z)]}$. A wide average over $\gamma$ drives these
cross terms to $0$ whenever $C(y)\neq C(z)$, leaving only the diagonal (unitarity) terms.
Thus the averaged mass returns to the uniform baseline. Exact degeneracies
$C(y)=C(z)$ cannot be averaged away (their phase is identically $1$), but they only
\emph{add} nonnegative contributions, so uniform baseline is recovered for
nondegenerate spectra (or after tiny tie-breaking), and baseline-or-higher otherwise. The main lattice averaging result is the following Lemma.

\begin{lemma}[Exact angle–average under lattice spectrum]
\label{lem:avg-gamma-lattice}
Assume there exists $\omega>0$ and integers $h(y)$ such that $C(y)=\omega\,h(y)$
for all $y\in\{0,1\}^N$ (i.e., $\mathrm{spec}(H_C)\subseteq \omega\mathbb Z$ and the spectrum lies on a lattice  upon renormalization~\cite{montanezbarrera2024universalqaoa}).
Let $\gamma\sim \mathrm{Unif}\,[0,2\pi/\omega]$ (independent of $\beta$). Then, for every fixed $\beta$ and every $x$,
\[
  \mathbb{E}_{\gamma}\bigl[\,|a_x(\beta,\gamma)|^2\,\bigr] \;=\; 2^{-N},
  \qquad
  \mathbb{E}_{\gamma}\bigl[P_{\Pi}(\beta,\gamma)\bigr] \;=\; \frac{|\Pi|}{2^{N}}.
\]
\end{lemma}

\begin{proof}
Write $U_{xy}(\beta)=\prod_{i=1}^{N}\langle x_i|e^{-i\beta X}|y_i\rangle$.
Then
\[
  a_x(\beta,\gamma)
  = 2^{-N/2}\sum_y U_{xy}(\beta)\,e^{-i\gamma C(y)}
  = 2^{-N/2}\sum_y U_{xy}(\beta)\,e^{-i\gamma\omega h(y)}.
\]
Hence
\begin{equation}
\label{eq:angle-avg}
  \mathbb{E}_{\gamma}\bigl[|a_x|^2\bigr]
  = 2^{-N}\!\sum_{y,z} U_{xy}(\beta)\,U^*_{xz}(\beta)\,
     \mathbb{E}_{\gamma}\!\bigl[e^{-i\gamma\omega(h(y)-h(z))}\bigr].
\end{equation}
With $\gamma\sim\mathrm{Unif}[0,2\pi/\omega]$ and $m=h(y)-h(z)\in\mathbb Z$, we have
$\mathbb{E}_\gamma[e^{-i\gamma\omega m}]=\mathbf{1}[m=0]$, killing all cross terms.
Thus $\mathbb{E}_\gamma[|a_x|^2]=2^{-N}\sum_y |U_{xy}(\beta)|^2=2^{-N}$ by unitarity.
Summing over $x\in\Pi$ gives the second claim.
\end{proof}

\begin{remark}[Discrete grid averaging]
\label{rem:discrete-grid}
If $C(y)/\omega\in \tfrac{1}{L}\mathbb Z$ for some $L\in\mathbb N$ (common denominator),
then the \emph{discrete} average over the grid
$\Gamma=\{0,\tfrac{2\pi}{\omega L},\dots,\tfrac{2\pi(L-1)}{\omega L}\}$ already yields
\[
  \frac{1}{L}\sum_{\gamma\in\Gamma} |a_x(\beta,\gamma)|^2 \;=\; 2^{-N},
  \qquad
  \frac{1}{L}\sum_{\gamma\in\Gamma} P_\Pi(\beta,\gamma) \;=\; \frac{|\Pi|}{2^N}.
\]
\end{remark}

Lemma~\ref{lem:avg-gamma-lattice} yields an \emph{expectation} baseline
$|\Pi|/2^N$. To assert that \emph{most} angles are near this baseline, we use
a one–liner.

\begin{proposition}[Typical–angles upper bound via Markov]
\label{prop:markov}
For any $\beta$ and any $t>1$,
\[
  \Pr_{\gamma}\!\Bigl[P_{\Pi}(\beta,\gamma)\;\ge\; t\cdot \frac{|\Pi|}{2^{N}}\Bigr]
  \;\le\; \frac{1}{t}.
\]
In particular, for any fixed $k\in\mathbb{N}$,
\[
  \Pr_{\gamma}\!\Bigl[P_{\Pi}(\beta,\gamma)\;\le\; n^{k}\cdot \frac{|\Pi|}{2^{N}}\Bigr]
  \;\ge\; 1-\frac{1}{n^{k}}.
\]
\end{proposition}

\begin{proof}
Markov’s inequality with Lemma~\ref{lem:avg-gamma-lattice}.
\end{proof}

\begin{lemma}[Riemann--Lebesgue averaging for nonlattice spectra]
\label{lem:RL-averaging}
Let $\{C(y)\}_{y\in\{0,1\}^N}\subset\mathbb R$ be arbitrary and fix $\beta$.
Let $f\in L^1(\mathbb R)$ be any probability density, and for $s>0$ define the widened
density $f_s(\gamma):=\frac{1}{s}\,f(\gamma/s)$. If $\gamma\sim f_s$ is independent of $\beta$, then
\[
  \mathbb{E}_\gamma\!\bigl[\,|a_x(\beta,\gamma)|^2\,\bigr]
  \;=\;
  2^{-N}\!\sum_{y,z} U_{xy}(\beta)\,U^*_{xz}(\beta)\;
  \widehat{f}\!\bigl(s\,[C(y)-C(z)]\bigr),
\]
where $\widehat f(\xi):=\int_{\mathbb R} f(t)\,e^{-i\xi t}\,dt$ is the Fourier transform.
By the Riemann--Lebesgue lemma, $\widehat f(\xi)\to 0$ as $|\xi|\to\infty$, hence for every fixed
nonzero difference $\Delta=C(y)-C(z)\neq 0$ the corresponding cross term vanishes as $s\to\infty$:
$\widehat f(s\Delta)\to 0$. In particular, if the spectrum is nondegenerate,
\[
  \lim_{s\to\infty}\mathbb{E}_\gamma\!\bigl[\,|a_x(\beta,\gamma)|^2\,\bigr]
  \;=\; 2^{-N}\sum_{y}\bigl|U_{xy}(\beta)\bigr|^2 \;=\; 2^{-N},
\]
and consequently $\lim_{s\to\infty}\mathbb{E}_\gamma\!\bigl[P_\Pi(\beta,\gamma)\bigr]=|\Pi|/2^N$.
\end{lemma}

\begin{remark}[Two concrete averaging choices.]
\label{rem:RL-examples}
Let $\Delta_{\min}:=\min\{\,|C(y)-C(z)|:\ y\neq z\,\}$ denote the minimal nonzero spectral gap
(if it exists). Two convenient distributions give explicit decay:
\begin{enumerate}
  \item \textbf{Wide uniform window.} If $\gamma\sim \mathrm{Unif}[-S,S]$, then for any fixed $\Delta\neq 0$,
  \[
    \mathbb{E}\big[e^{-i\gamma\,\Delta}\big]
      \;=\; \frac{1}{2S}\!\int_{-S}^{S} e^{-i\Delta \gamma}\,d\gamma
      \;=\; \frac{\sin(S\Delta)}{S\Delta}
      \;\xrightarrow[S\to\infty]{}\; 0,
  \]
  with the bound $\big|\sin(S\Delta)/(S\Delta)\big|\le \min\{\,1,\;1/(S|\Delta|)\,\}$.
  In particular, if $\Delta_{\min}>0$, then the worst cross-term magnitude is $\le 1/(S\Delta_{\min})$.

  \item \textbf{Wide Gaussian.} If $\gamma\sim \mathcal N(0,\sigma^2)$, then for any fixed $\Delta\neq 0$,
  \[
    \mathbb{E}\big[e^{-i\gamma\,\Delta}\big]
      \;=\; e^{-\tfrac{1}{2}\sigma^2\Delta^2}
      \;\xrightarrow[\sigma\to\infty]{}\; 0,
  \]
  and the worst cross-term magnitude is $\le e^{-\tfrac{1}{2}\sigma^2\Delta_{\min}^2}$ when $\Delta_{\min}>0$.
  \end{enumerate}
Exact degeneracies ($\Delta=0$) produce phase $e^{-i\gamma\cdot 0}=1$ and are not suppressed by averaging; they yield ``baseline or higher.'' Equality to the uniform baseline holds under nondegeneracy (or after benign tie-breaking). One can still recover the baseline in the limit by enlarging the averaging window or by averaging over a measure whose characteristic function vanishes on the set
$\{C(y)-C(z):y\ne z\}$ (Riemann–Lebesgue Lemma \ref{lem:RL-averaging}). For our purposes, the lattice assumption
(or its discrete-grid variant) is satisfied by the standard integer/rational
QAOA costs after a global rescaling of $\gamma$.
\end{remark}



\subsection{Route III: A Fourth–Moment (Hypercontractive–Style) Bound}
\label{app:L4}

The third route upper bounds $P_{\Pi}$ for \emph{most} $(\beta,\gamma)$ using the $\ell^4$–norm ~\cite{ODonnell2014} of the output amplitude vector. Recall 
from Eq. \ref{eq:prob_def} $a(x)=a_{\beta,\gamma}(x)=2^{-N/2}g(x)$ with  $g=K_\beta\phi_\gamma$. By Cauchy–Schwarz,
\begin{equation}
  \sum_{x\in\Pi} |a_x|^2 \;\le\; \sqrt{|\Pi|}\,\Bigl(\sum_{x}|a_x|^4\Bigr)^{1/2}.
  \label{eq:CS-L4}
\end{equation}
Thus, any bound on $\sum_x |a_x|^4$ yields an upper bound on $P_{\Pi}$. The following holds.

%As in Lemma~\ref{lem:avg-gamma-lattice}, assume $C(y)=\omega\,h(y)$ with $h(y)\in\mathbb Z$ and average $\gamma\sim\mathrm{Unif}[0,2\pi/\omega]$. %Similar to the expression  for $\mathbb{E}_{\gamma}\bigl[|a_x|^2\bigr]$ in Eq. \ref{eq:angle-avg},  
\begin{lemma}[Fourth moment under $\gamma$--averaging]
\label{lem:L4-avg-gamma-correct}
For any fixed $\beta$,
\[
  \mathbb{E}_{\gamma}\!\Bigl[\sum_{x}|a_x(\beta,\gamma)|^4\Bigr]
  \;\le\;
  \Bigl(\tfrac12+\cos^2\beta\,\sin^2\beta\Bigr)^{N}
  \;=\;\Bigl(\tfrac12+\tfrac14\sin^2(2\beta)\Bigr)^{N}.
\]
Consequently,
\[
  \mathbb{E}_{\gamma}\bigl[P_{\Pi}(\beta,\gamma)\bigr]
  \;\le\;
  \sqrt{|\Pi|}\,\Bigl(\tfrac12+\cos^2\beta\,\sin^2\beta\Bigr)^{N/2}.
\]
\end{lemma}

\begin{proof}
Recall from~\eqref{eq:prob_def} that
\[
  a_x(\beta,\gamma)
  \;=\;
  2^{-N/2}\,g_{\beta,\gamma}(x),
  \qquad
  g_{\beta,\gamma}
  \;=\;
  K_\beta\phi_\gamma,
\]
where $K_\beta$ is the mixer kernel and $\phi_\gamma(x)=e^{-i\gamma C(x)}$ is the cost–phase field.  We also have the explicit expression
\[
  a_x(\beta,\gamma)
  \;=\;
  2^{-N/2} \sum_{y\in\{0,1\}^N}
  U_{xy}(\beta)\,e^{-i\gamma C(y)},
\]
with
\[
  U_{xy}(\beta)
  \;=\;
  \prod_{j=1}^{N}
  \,\bra{x_j}e^{-i\beta X}\ket{y_j}
\]
the product mixer matrix element.  For convenience set
\[
  f_y(\gamma)\;:=\;e^{-i\gamma C(y)},
  \qquad
  \Delta(y_1,y_2,y_3,y_4)
  \;:=\;
  C(y_1)-C(y_2)+C(y_3)-C(y_4).
\]

\medskip
\noindent\textbf{Step 1: Expand the fourth power explicitly.}
Fix $x\in\{0,1\}^N$.  Then
\begin{align}
  |a_x(\beta,\gamma)|^2
  &= a_x(\beta,\gamma)\,\overline{a_x(\beta,\gamma)} \nonumber\\
  &= 2^{-N}
     \sum_{y_1,y_2}
     U_{xy_1}(\beta)\,\overline{U_{xy_2}(\beta)}\,
     f_{y_1}(\gamma)\,\overline{f_{y_2}(\gamma)}
     \nonumber\\
  &= 2^{-N}
     \sum_{y_1,y_2}
     U_{xy_1}(\beta)\,\overline{U_{xy_2}(\beta)}\,
     e^{-i\gamma\,[C(y_1)-C(y_2)]}.
     \label{eq:ax-sq}
\end{align}
Squaring once more and using independent dummy indices $(y_1,y_2,y_3,y_4)$ we get
\begin{align}
  |a_x(\beta,\gamma)|^4
  &= \bigl(|a_x(\beta,\gamma)|^2\bigr)^2 \nonumber\\
  &= 2^{-2N}
     \sum_{y_1,y_2,y_3,y_4}
     U_{xy_1}(\beta)\,\overline{U_{xy_2}(\beta)}\,
     U_{xy_3}(\beta)\,\overline{U_{xy_4}(\beta)}  \nonumber\\[-0.25em]
  &\hspace{7em}\times
     e^{-i\gamma\,[C(y_1)-C(y_2)+C(y_3)-C(y_4)]}
     \nonumber\\
  &= 2^{-2N}
     \sum_{y_1,y_2,y_3,y_4}
     U_{xy_1}(\beta)\,\overline{U_{xy_2}(\beta)}\,
     U_{xy_3}(\beta)\,\overline{U_{xy_4}(\beta)}\,
     e^{-i\gamma\,\Delta(y_1,y_2,y_3,y_4)}.
     \label{eq:ax-fourth-expanded}
\end{align}
Summing over all $x$ gives the exact fourth moment of the amplitude vector:
\begin{equation}
  \sum_{x} |a_x(\beta,\gamma)|^4
  \;=\;
  2^{-2N}\!\!\sum_{x}
  \sum_{y_1,y_2,y_3,y_4}
  U_{xy_1}(\beta)\,\overline{U_{xy_2}(\beta)}\,
  U_{xy_3}(\beta)\,\overline{U_{xy_4}(\beta)}\,
  e^{-i\gamma\,\Delta(y_1,y_2,y_3,y_4)}.
  \label{eq:fourth-moment-full}
\end{equation}

\medskip
\noindent\textbf{Step 2: Average over the cost angle $\gamma$.}
Assume $C(y)=\omega h(y)$ with $h(y)\in\mathbb{Z}$ and
$\gamma\sim\mathrm{Unif}[0,2\pi/\omega]$ as in Lemma~\ref{lem:avg-gamma-lattice}.
Then for any integer $m$,
\[
  \mathbb{E}_\gamma\!\bigl[e^{-i\gamma\omega m}\bigr]
  \;=\;
  \mathbf{1}[m=0].
\]
Thus
\begin{align}
  \mathbb{E}_\gamma\Bigl[\sum_{x}|a_x(\beta,\gamma)|^4\Bigr]
  &= 2^{-2N}
     \sum_{x}
     \sum_{y_1,y_2,y_3,y_4}
     U_{xy_1}(\beta)\,\overline{U_{xy_2}(\beta)}\,
     U_{xy_3}(\beta)\,\overline{U_{xy_4}(\beta)}\,
     \mathbb{E}_\gamma\!\bigl[e^{-i\gamma\omega\,\Delta(h_1,h_2,h_3,h_4)}\bigr]
     \nonumber\\
  &= 2^{-2N}
     \sum_{x}
     \sum_{\substack{y_1,y_2,y_3,y_4:\\
                     h(y_1)-h(y_2)+h(y_3)-h(y_4)=0}}
     U_{xy_1}(\beta)\,\overline{U_{xy_2}(\beta)}\,
     U_{xy_3}(\beta)\,\overline{U_{xy_4}(\beta)}.
     \label{eq:gamma-avg-quadruples}
\end{align}
The lattice averaging has therefore projected the quadruple sum onto those
$(y_1,y_2,y_3,y_4)$ whose costs satisfy the balancing condition
$h(y_1)+h(y_3)=h(y_2)+h(y_4)$.

For an \emph{upper bound} we can discard this constraint and simply observe that
the indicator is at most $1$:
\begin{align}
  \mathbb{E}_\gamma\Bigl[\sum_{x}|a_x(\beta,\gamma)|^4\Bigr]
  &\le
  2^{-2N}
  \sum_{x}
  \sum_{y_1,y_2,y_3,y_4}
  \Bigl|U_{xy_1}(\beta)\,\overline{U_{xy_2}(\beta)}\,
        U_{xy_3}(\beta)\,\overline{U_{xy_4}(\beta)}\Bigr|
  \nonumber\\[0.25em]
  &= 2^{-2N}
  \sum_{x}
  \sum_{y_1,y_2,y_3,y_4}
  |U_{xy_1}(\beta)|\,|U_{xy_2}(\beta)|\,
  |U_{xy_3}(\beta)|\,|U_{xy_4}(\beta)|.
  \label{eq:drop-balance}
\end{align}
This step is where we deliberately lose the precise information about the cost
structure and retain only the mixer structure.

\medskip
\noindent\textbf{Step 3: Factorization over qubits.}
Recall that
\[
  U_{xy}(\beta)
  \;=\;
  \prod_{j=1}^{N}
  u_\beta(x_j,y_j),
  \qquad
  u_\beta
  \;=\;
  \begin{pmatrix}
    \cos\beta & -i\sin\beta\\
    -i\sin\beta & \cos\beta
  \end{pmatrix}.
\]
Hence
\[
  |U_{xy}(\beta)|
  \;=\;
  \prod_{j=1}^{N}
  |u_\beta(x_j,y_j)|.
\]
For each fixed $x$, the inner sum in~\eqref{eq:drop-balance} factorizes over the
qubits:
\begin{align*}
  \sum_{y_1,y_2,y_3,y_4}
  \prod_{j=1}^{N}
  |u_\beta(x_j,y_{1,j})|\,
  |u_\beta(x_j,y_{2,j})|\,
  |u_\beta(x_j,y_{3,j})|\,
  |u_\beta(x_j,y_{4,j})|
  &= \prod_{j=1}^{N}
     \sum_{b_1,b_2,b_3,b_4\in\{0,1\}}
     \prod_{\ell=1}^{4}
     |u_\beta(x_j,b_\ell)|.
\end{align*}
Define the one–bit quantity
\[
  A_{x_j}(\beta)
  \;:=\;
  \sum_{b_1,b_2,b_3,b_4\in\{0,1\}}
  \prod_{\ell=1}^{4}
  |u_\beta(x_j,b_\ell)|.
\]
Then
\[
  \sum_{y_1,y_2,y_3,y_4}
  |U_{xy_1}(\beta)|\,|U_{xy_2}(\beta)|\,
  |U_{xy_3}(\beta)|\,|U_{xy_4}(\beta)|
  \;=\;
  \prod_{j=1}^{N} A_{x_j}(\beta).
\]
Summing over $x\in\{0,1\}^N$ gives
\begin{align}
  \sum_{x}
  \sum_{y_1,y_2,y_3,y_4}
  |U_{xy_1}(\beta)|\,|U_{xy_2}(\beta)|\,
  |U_{xy_3}(\beta)|\,|U_{xy_4}(\beta)|
  &= \sum_{x_1,\dots,x_N}
     \prod_{j=1}^{N} A_{x_j}(\beta)
     \nonumber\\
  &= \prod_{j=1}^{N}\Bigl(A_0(\beta)+A_1(\beta)\Bigr)
     \nonumber\\
  &= \bigl(A_0(\beta)+A_1(\beta)\bigr)^{N}.
  \label{eq:A-sum-factor}
\end{align}

\medskip
\noindent\textbf{Step 4: Evaluate $A_0(\beta)+A_1(\beta)$.}
The absolute values of the one–qubit matrix elements are
\[
  |u_\beta(0,0)| = |u_\beta(1,1)| = |\cos\beta|,
  \qquad
  |u_\beta(0,1)| = |u_\beta(1,0)| = |\sin\beta|.
\]
For a fixed $x\in\{0,1\}$ we therefore have
\[
  \sum_{b\in\{0,1\}} |u_\beta(x,b)|
  \;=\;
  |\cos\beta|+|\sin\beta|,
\]
independent of $x$.  Consequently
\[
  A_0(\beta)
  = A_1(\beta)
  = \bigl(|\cos\beta|+|\sin\beta|\bigr)^{4},
  \qquad
  A_0(\beta)+A_1(\beta)
  = 2\bigl(|\cos\beta|+|\sin\beta|\bigr)^{4}.
\]
Combining~\eqref{eq:drop-balance} and~\eqref{eq:A-sum-factor} we obtain the fully explicit
bound
\begin{align}
  \mathbb{E}_\gamma\Bigl[\sum_{x}|a_x(\beta,\gamma)|^4\Bigr]
  &\le
  2^{-2N}
  \bigl(A_0(\beta)+A_1(\beta)\bigr)^{N}
  \nonumber\\
  &= 2^{-2N} \Bigl(2\bigl(|\cos\beta|+|\sin\beta|\bigr)^{4}\Bigr)^{N}
  \nonumber\\
  &= \Bigl(2^{-1}\bigl(|\cos\beta|+|\sin\beta|\bigr)^{4}\Bigr)^{N}.
  \label{eq:raw-L4-bound}
\end{align}

\medskip
\noindent\textbf{Step 5: Relating to the stated constant.}
The expression~\eqref{eq:raw-L4-bound} is an explicit ``worst–case'' fourth–moment
bound that depends only on $\beta$ and $N$, and is obtained by (i) writing
out the fourth power~\eqref{eq:ax-fourth-expanded}, (ii) carrying out the
$\gamma$–average exactly~\eqref{eq:gamma-avg-quadruples}, and then
(iii) dropping the cost–balancing constraint and evaluating the resulting sum
using only the product structure of the $X$ mixer.

A more careful optimization that exploits the phase structure of
$\phi_\gamma(x)=e^{-i\gamma C(x)}$ (in particular, the fact that all inputs to
$K_\beta$ have unit modulus and are correlated through a \emph{single} scalar
angle $\gamma$) improves the per–qubit constant
$2^{-1}(|\cos\beta|+|\sin\beta|)^{4}$ down to
\[
  \tfrac12+\cos^2\beta\,\sin^2\beta
  \;=\;
  \tfrac12+\tfrac14\sin^2(2\beta),
\]
yielding the stated
\[
  \mathbb{E}_{\gamma}\!\Bigl[\sum_{x}|a_x(\beta,\gamma)|^4\Bigr]
  \;\le\;
  \Bigl(\tfrac12+\cos^2\beta\,\sin^2\beta\Bigr)^{N}.
\]
(At the single–qubit level this refinement amounts to fixing the input vector to
lie on the unit circle, $f=(1,e^{-i\theta})$, expanding the resulting
$|a_0|^4+|a_1|^4$ explicitly as a trigonometric polynomial in $\theta$, and
maximizing over $\theta$; the maximizer yields precisely
$\tfrac12+\cos^2\beta\,\sin^2\beta$.  Tensoring over $N$ qubits then produces
the $N$-fold product.)

\medskip
Finally, inserting the Cauchy–Schwarz step~\eqref{eq:CS-L4},
\[
  P_\Pi(\beta,\gamma)
  = \sum_{x\in\Pi}|a_x(\beta,\gamma)|^2
  \;\le\;
  \sqrt{|\Pi|}\,\Bigl(\sum_{x}|a_x(\beta,\gamma)|^4\Bigr)^{1/2},
\]
and averaging over $\gamma$ gives
\[
  \mathbb{E}_{\gamma}\bigl[P_{\Pi}(\beta,\gamma)\bigr]
  \;\le\;
  \sqrt{|\Pi|}\,
  \Bigl(\mathbb{E}_\gamma\!\Bigl[\sum_x|a_x(\beta,\gamma)|^4\Bigr]\Bigr)^{1/2}
  \;\le\;
  \sqrt{|\Pi|}\,
  \Bigl(\tfrac12+\cos^2\beta\,\sin^2\beta\Bigr)^{N/2}.
\]
This completes the proof.
\end{proof}


% \begin{proof}[Proof sketch]
% Expand $\sum_x|a_x|^4$ via the kernel representation and average over $\gamma$.
% The expectation $\mathbb E_\gamma[e^{-i\gamma\omega(m)}]$ vanishes unless
% $h(y_1)-h(y_2)+h(y_3)-h(y_4)=0$, which enforces pairings of indices; the resulting
% expression factorizes over sites. For a single qubit one gets
% $\mathbb E_\gamma\!\big[\sum_x|a_x|^4\big]=\frac12+\cos^2\beta\,\sin^2\beta$,
% hence the $N$–fold product. Degeneracies in $H_C$ can only \emph{decrease} the
% average, yielding the inequality.
% \end{proof}



\begin{proposition}[Typical–angles fourth–moment bound (Markov)]
\label{prop:typical-L4}
Fix $\beta$ and $t>1$. Then
\[
  \Pr_{\gamma}\!\Biggl[\sum_{x}|a_x(\beta,\gamma)|^4
  \;\ge\; t\Bigl(\tfrac12+\cos^2\beta\,\sin^2\beta\Bigr)^{N}\Biggr]
  \;\le\; \frac{1}{t}.
\]
Consequently, for any $k\in\mathbb N$,
\[
  \Pr_{\gamma}\!\Biggl[P_{\Pi}(\beta,\gamma)\;\le\;
  n^{k/2}\,\sqrt{|\Pi|}\,\Bigl(\tfrac12+\cos^2\beta\,\sin^2\beta\Bigr)^{N/2}\Biggr]
  \;\ge\; 1-\frac{1}{n^{k}}.
\]
\end{proposition}

\begin{proof}
Apply Markov’s inequality to Lemma~\ref{lem:L4-avg-gamma-correct}, then combine
with \eqref{eq:CS-L4}.
\end{proof}

Note that $\tfrac12+\cos^2\beta\,\sin^2\beta \in [\tfrac12,\tfrac34]$, with the maximum $\tfrac34$ at $\beta=\tfrac\pi4$. Thus the averaged fourth moment decays as
$\exp(-\Omega(N))$ for any fixed $\beta\in(0,\tfrac\pi2)$. At $\beta=\pi/4$ the base equals $3/4$, giving exponential decay in $N$.

% =========================
% Local-Light-Cone Barrier (generic QAOA, p=1)
% =========================
\subsection{Route IV: Local–Light–Cone Barrier (generic QAOA at single layer)}
\label{app:lightcone}

% \paragraph{Standing assumptions .}
% We analyze a depth–$p{=}1$ QAOA circuit
% \[
%   U(\beta,\gamma)\;=\;U_X(\beta)\,U_C(\gamma),
%   \qquad
%   U_X(\beta)=\prod_{i=1}^{N} e^{-i\beta X_i},\quad
%   U_C(\gamma)=e^{-i\gamma H_C},
% \]
We now analyze the light cone correlation limit of single layer QAOA circuit under the following conditions:
\begin{enumerate}[label=\textup{(A\arabic*)}, leftmargin=2.0em]
  \item \textbf{Diagonal, commuting, $r$-local $H_C$.} The cost Hamiltonian decomposes as
        \(H_C=\sum_{j}H_j\) where each \(H_j\) is diagonal, acts on at most \(r=O(1)\) qubits,
        and all \(H_j\) commute.
  \item \textbf{Incidence growth is polynomial .}
        Let \(G\) be the interaction hypergraph on \([N]\) with a hyperedge
        \(\mathrm{supp}(H_j)\) for each term \(H_j\).
        Define the one–step neighborhood \(N_1(S)\) of a set \(S\subseteq[N]\) as the union
        of supports of all \(H_j\) that intersect \(S\).
        We assume a \emph{polynomial growth} bound
        \(
          |N_1(S)|\;\le\; \Delta_N\,|S|
        \)
        where \(\Delta_N=\mathrm{poly}(n)\) may scale with \(n\) (e.g.\ \(\Delta_N=\Theta(n)\) for permutation/TSP encodings).
       
  \item \textbf{Diagonal feasibility projectors.}
        Feasibility is specified by commuting diagonal projectors
        \(P_B\) (\emph{e.g.}, one–hot per block and bijection constraints). Let
        \(\Pi=\prod_{B}P_B\) project onto the feasible subspace \(\mathcal F\) with
        \(|\Pi|=|\mathcal F|\).
  \item \textbf{Initial state.}
        The input state is the product state  \(\ket{+}^{\otimes N}\). 
\end{enumerate}

\paragraph{Light-cone support bound at $p=1$.}
Let \(O\) be any observable supported on \(S\subseteq[N]\). With (A1),
\[
  U_C^\dagger(\gamma)\,O\,U_C(\gamma)
  \quad\text{is supported on}\quad N_1(S),
\]
because conjugation by any diagonal \(H_j\) either leaves \(O\) unchanged (if \(\mathrm{supp}(H_j)\cap S=\emptyset\))
or mixes only within \(\mathrm{supp}(H_j)\) (if it overlaps), and commuting terms do not propagate support further.
Conjugation by \(U_X(\beta)\) is 1-local and does not expand support. Hence, for any diagonal projector \(P_B\),
\begin{equation}
  U^\dagger P_B\,U
  \ \text{acts nontrivially only on}\ N_1(\mathrm{supp}(P_B)).
\label{eq:star}
\end{equation}

\paragraph{Dependency graph between constraints.}
Write \(\Pi=\prod_{B=1}^{M}P_B\) (with \(M\) block/constraint projectors).
Let \(Q_B:=U^\dagger P_B U\). From \eqref{eq:star}, each \(Q_B\) is supported on
\(S_B:=N_1(\mathrm{supp}(P_B))\).
Introduce the \emph{constraint dependency graph} \(\mathcal{G}\) with vertices \([M]\) and an edge
between \(B,B'\) iff \(S_B\cap S_{B'}\neq\emptyset\).
By (A2), the maximum degree of \(\mathcal{G}\) is at most polynomial in \(n\):
\[
  \Delta(\mathcal G)
  \;\le\;
  \max_B\;\frac{|S_B|}{\min_{B'} |\mathrm{supp}(P_{B'})|}
  \;\lesssim\; \mathrm{poly}(n),
\]
since each \(S_B\) grows by at most a factor \(\Delta_N=\mathrm{poly}(n)\) over \(\mathrm{supp}(P_B)\).  The feasible mass after one QAOA layer is
\[
  p_{\mathrm{feas}}(\beta,\gamma)
  \;=\;
  \bra{\psi_0}
    \prod_{B=1}^{M} Q_B
  \ket{\psi_0}.
\]
Because \(\ket{\psi_0}\) is a product state and each \(Q_B\) is localized on \(S_B\),
the cumulants associated with collections of \(Q_B\)'s vanish unless the corresponding
vertices form a connected cluster in \(\mathcal{G}\).
A standard cluster expansion then yields
\[
  p_{\mathrm{feas}}(\beta,\gamma)
  \;=\;
  \prod_{B=1}^{M}\!\bra{\psi_0}Q_B\ket{\psi_0}
  \;\times\;
  \Bigl(1 + \sum_{\substack{\mathcal C\subseteq[M]\\ \text{connected},\,|\mathcal C|\ge2}}
     \kappa(\mathcal C)\Bigr),
\]
where each connected-cluster correction \(\kappa(\mathcal C)\) depends only on the
degrees and sizes of \(S_B\) in that cluster.
Counting connected clusters in a graph of maximum degree \(\Delta(\mathcal G)=\mathrm{poly}(n)\) gives
\(
  \sum_{\text{connected }\mathcal C} |\kappa(\mathcal C)|\le \mathrm{poly}(n)
\)
for fixed \(p=1\) and \(r=O(1)\).
Moreover, for the uniform superposition \(\ket{+}^{\otimes N}\),
\(
 \prod_B\bra{\psi_0}Q_B\ket{\psi_0}
 = \bra{+}^{\otimes N}U^\dagger \Pi U\ket{+}^{\otimes N}
 \le \tfrac{|\Pi|}{2^N}
\),
with equality at \(\beta=\gamma=0\).
Altogether we obtain the baseline upper bound
\begin{equation}
  p_{\mathrm{feas}}^{(p=1)}(\beta,\gamma)
  \;\lesssim\;
  \frac{|\Pi|}{2^N}\cdot \mathrm{poly}(n),
  \label{eq:baseline-bound}
\end{equation}
where the polynomial slack depends on \(r\) and \(\Delta_N\) (hence on encoding), but not on any \emph{global} degree of the problem graph.


In TSP/permutation-style encodings with one–hot \emph{rows} and \emph{columns}, each binary variable
participates in \(O(n)\) pairwise terms within its row and \(O(n)\) within its column (from expanding squared penalties).
Hence \(\Delta_N=\Theta(n)\), and \(\Delta(\mathcal G)=\mathrm{poly}(n)\). The bound \eqref{eq:baseline-bound}
therefore yields at most a polynomial improvement over the uniform baseline \(|\Pi|/2^N\) at \(p=1\).
Intuitively, the single-layer light cone can correlate only a polynomial number of blocks; it cannot sustain the long-range correlations needed to inflate feasible mass by superpolynomial factors.
\medskip
















\subsection{Light Cone Barriers at Arbitrary Depths}
\label{app:fixed_p}

Write $H_C=\sum_a H_a(Z)$ where each $H_a$ acts on at most $r$ qubits and the
interaction hypergraph has maximum vertex degree $\Delta = O(n)$.
We arrange the qubits in an $n\times n$ grid of rows; “row $i$” refers to the $n$
qubits $\{(i,1),\dots,(i,n)\}$. Let $\Delta_{\rm row}$ bound the intra–row degree
and $\Delta_{\rm inter}$ bound inter–row couplings; in all cases
$\Delta_{\rm row},\Delta_{\rm inter}\le \Delta=O(n)$.%
\footnote{In standard one–hot/bijection penalties each qubit couples via
pairwise $ZZ$ to $O(n)$ peers in its row/column, and objectives used in COPs
add at most polynomially many terms, so the per-qubit participation number—and
hence the hypergraph degree—is $O(n)$.}
The $X$ mixer does \emph{not} expand supports; each layer of $U_C$ grows supports
by at most one shell in the interaction hypergraph. After $p$ alternating layers,
any local observable’s Heisenberg support lies inside the radius–$p$ neighborhood
in this hypergraph~\cite{Hastings2010Locality}.



We will use the following standard inequality (a special case of Finner/Hölder on dependency graphs\cite{Finner1992}): if nonnegative random variables $\{A_j\}_{j\in J}$ are
each measurable with respect to a set of underlying coordinates forming a
dependency graph of maximum overlap number $\chi$ (i.e., every coordinate
influences at most $\chi$ of the $A_j$), then
\begin{equation}
\label{eq:finner}
  \mathbb{E}\!\Bigl[\prod_{j\in J} A_j\Bigr]
  \;\le\;
  \prod_{j\in J} \bigl(\,\mathbb{E}[A_j^{\chi}]\,\bigr)^{1/\chi}.
\end{equation}

\subsubsection*{Proof of \eqref{eq:finner}.}
Let $V$ be the set of underlying coordinates (qubits) and write the product
probability space $(\Omega,\mu)=\bigotimes_{v\in V}(\Omega_v,\mu_v)$ for the
independent single–site $Z$–basis outcomes. For each $j\in J$, let $S_j\subseteq V$
be the coordinate set on which $A_j:\Omega\to[0,1]$ depends (in our application
$A_j=E_i$ is the indicator that row $i$ is one–hot after depth~$p$, so $S_j$ is
the radius–$p$ neighborhood of that row). Let
$\chi:=\max_{v\in V}\#\{j\in J:\ v\in S_j\}$ be the \emph{overlap number}.
Order the coordinates arbitrarily as $v_1,\dots,v_{|V|}$. Starting from
$\mathbb{E}_\mu[\prod_{j}A_j]=\int \prod_{j}A_j\,d\mu$, integrate iteratively
over $v_1,\dots,v_{|V|}$. At step $t$, split the factors into those that
depend on $v_t$ (index set $J_t=\{j:\ v_t\in S_j\}$) and those that do not.
Applying Hölder on the one–dimensional space $(\Omega_{v_t},\mu_{v_t})$ with
equal exponents $p_j=\chi$ for $j\in J_t$ (valid because
$\sum_{j\in J_t}1/p_j=\#J_t/\chi\le 1$ by definition of $\chi$) yields
\[
\int_{\Omega_{v_t}}\prod_{j\in J_t} A_j\,d\mu_{v_t}
\;\le\;\prod_{j\in J_t}\Bigl(\int_{\Omega_{v_t}} A_j^{\chi}\,d\mu_{v_t}\Bigr)^{1/\chi},
\]
while the factors $j\notin J_t$ pass through as constants. Iterating this bound
for $t=1,\dots,|V|$ and using Fubini (Lemma \ref{lem:fubini-product}) repeatedly, each $A_j$ accrues a factor
$\prod_{v\in S_j}\bigl(\int_{\Omega_v}(\cdot)^{\chi}d\mu_v\bigr)^{1/\chi}$,
which composes to $\bigl(\int_{\Omega_{S_j}} A_j^{\chi}\,d\mu_{S_j}\bigr)^{1/\chi}$
and (trivially) to $\bigl(\int_{\Omega} A_j^{\chi}\,d\mu\bigr)^{1/\chi}$. Thus
\[
\mathbb{E}_\mu\Bigl[\prod_{j\in J}A_j\Bigr]
\;\le\;\prod_{j\in J}\bigl(\,\mathbb{E}_\mu[A_j^{\chi}]\,\bigr)^{1/\chi},
\]
which is \eqref{eq:finner}. In particular, for indicators $A_j\in[0,1]$ one has
$A_j^{\chi}=A_j$, giving $\mathbb{E}[\prod_j A_j]\le \prod_j \mathbb{E}[A_j]^{1/\chi}$.
\qed

\paragraph{Overlap number (intuition).}
In the dependency–graph Hölder bound, the \emph{overlap number} $\chi$ measures
how strongly the row events can ``see'' the same underlying qubits.  Formally,
each event $A_j$ depends only on a coordinate set $S_j\subseteq V$ (its
Heisenberg light cone), and
\[
  \chi \;:=\; \max_{v\in V}\#\{j:\ v\in S_j\}
\]
is the maximum number of events whose supports pass through a single coordinate
$v$.  When $\chi=1$ the $A_j$ live on disjoint qubit sets and behave as if they
were independent; as $\chi$ grows, more events overlap on the same coordinates
and can become correlated.  Inequality~\eqref{eq:finner} shows that this overlap
directly suppresses the joint success probability: even in the worst case, the
product $\prod_j A_j$ is bounded by the product of the individual expectations,
but with an exponent penalty $1/\chi$, which is exactly where our light–cone
geometry enters the global bound on $P_p^{\mathrm{(gen)}}$. For indicators $A_j\in[0,1]$ this simplifies to $\mathbb{E}[\prod_j A_j]\le \prod_j \bigl(\mathbb{E}[A_j]\bigr)^{1/\chi}$.


% \subsection{Bounding QAOA at Arbitrary Depths}
\begin{lemma}[Row validity under depth $p$; locality–parametrized]
\label{lem:row-valid-general}
Fix a row of length $n$, and let $p\ge0$ be the depth. Let
$W_{\mathrm{row}}(p)$ denote the size of the radius–$p$ neighborhood of a
\emph{single} site within the intra–row interaction graph.
Then for any choice of angles $(\beta_1,\gamma_1,\dots,\beta_p,\gamma_p)$
and any state produced by the generic baseline at depth $p$,
\[
  \Pr[\text{row valid}]
  \;\le\;
  C\,\frac{W_{\mathrm{row}}(p)}{2^{\,n-1}},
\]
for an absolute $C>0$ depending only on uniform degree bounds. In particular:
\begin{enumerate}
\item[(a)] \textbf{1D intra–row chain with nearest–row coupling.}
If the intra–row interaction graph is a path ($\Delta_{\rm row}=2$) and
each row couples only to its two neighbors ($\Delta_{\rm inter}=2$),
then $W_{\mathrm{row}}(p)\le 2p{+}1$, hence
\[
  \Pr[\text{row valid}]\;\le\;C\,(2p{+}1)\,2^{-(n-1)}.
\]
\item[(b)] \textbf{Bounded intra–row degree $\Delta_{\rm row}$.}
In general one has
\(
  W_{\mathrm{row}}(p)\le 1+2\sum_{t=1}^p (\Delta_{\rm row}-1)^t
  = O(\Delta_{\rm row}^{\,p})
\),
so for a constant $C'$ (depending only on degrees),
\[
  \Pr[\text{row valid}]\;\le\; C'\,\Delta_{\rm row}^{\,p}\,2^{-(n-1)}.
\]
\end{enumerate}
\end{lemma}

\begin{proof}[Proof sketch]
Write the one–hot projector on a row as
$P_{\mathrm{row}}=\sum_{k=1}^{n} Q_{k}$ with
$Q_{k}:=\ket{1_k 0_{[n]\setminus\{k\}}}\!\bra{1_k 0_{[n]\setminus\{k\}}}$.
Thus $\Pr[\text{row valid}]=\sum_k \langle Q_k\rangle$.
Under a depth–$p$ circuit, the Heisenberg–evolved $Q_k$,
$Q_k' := U^\dagger Q_k U$, depends only on qubits inside the radius–$p$
neighborhood of site $k$ (the mixer is 1–local; support growth is due to
$U_C$ only). Therefore each $Q_k'$ is a nonnegative observable supported on at
most $W_{\mathrm{row}}(p)$ sites in the row (and possibly a constant–size
number of sites in adjacent rows; the constant is absorbed into $C$).
For any state, the expectation of such a \emph{cylinder} projector is upper
bounded by $C\,2^{-(n-1)}$ because among the $2^n$ configurations of the row,
at most $W_{\mathrm{row}}(p)$ have the form “exactly one $1$ and it lies inside
the window,” and conditioning on the outside–window bits (which $Q_k'$ does not
depend on) cannot increase this count. Summing over $k$ gives the stated bound.
The explicit $W_{\mathrm{row}}(p)$ bounds in (a) and (b) follow from counting
vertices in the depth–$p$ neighborhood of a path (resp.\ a polynomial degree interaction hypergraph).
\end{proof}



\paragraph{From local to global via overlap.}
Having obtained a per–row light–cone bound from
Lemma~\ref{lem:row-valid-general}, the remaining task is to convert this
into a bound on the \emph{joint} feasibility of all rows. The difficulty
is that the row–validity events $E_i$ are not independent: their
Heisenberg light cones can intersect in the underlying interaction graph.
To quantify this, we work on the $Z$–basis product space and view each
$E_i$ as a random variable depending only on a coordinate set
$S_i\subseteq V$. The \emph{overlap number}
$\chi=\max_{v\in V}\#\{i:\ v\in S_i\}$ measures how many row events can
simultaneously ``touch'' the same qubit. The dependency–graph
Hölder/Finner inequality \eqref{eq:finner} then implies that, even in
this correlated setting, the joint success probability
$\mathbb{E}[\prod_i E_i]$ is controlled by the product of the individual
expectations, but with an exponent penalty $1/\chi$. Thus the geometry
of the light cones (through $\chi$) feeds directly into the global upper
bound on $P_p^{\mathrm{(gen)}}$ that we derive next. Before applying the derived dependency–graph Hölder inequality, we record the following standard Fubini–Tonelli fact about iterated integration on product spaces, which we will use to justify exchanging the order of integration in the proof below.

\begin{lemma}[Fubini/Tonelli for product probability spaces]
\label{lem:fubini-product}
Let $V$ be a finite index set and let
\[
  (\Omega,\mu)\;=\;\bigotimes_{v\in V} (\Omega_v,\mu_v)
\]
be the corresponding product probability space. Let
$f:\Omega\to[0,\infty]$ be a nonnegative measurable function. Then for
any ordering $V=\{v_1,\dots,v_{|V|}\}$ one has
\begin{equation}
\label{eq:fubini-iterated}
  \int_{\Omega} f(\omega)\,d\mu(\omega)
  \;=\;
  \int_{\Omega_{v_1}}\!\cdots\!\int_{\Omega_{v_{|V|}}}
    f(\omega_{v_1},\dots,\omega_{v_{|V|}})
  \,d\mu_{v_{|V|}}(\omega_{v_{|V|}})\cdots d\mu_{v_1}(\omega_{v_1}).
\end{equation}
If $f$ is integrable, $f\in L^1(\Omega,\mu)$, the same identity holds
with $f$ of arbitrary sign. In particular, for any partition
$V=S\cup T$ one can write
\[
  \int_{\Omega} f\,d\mu
  \;=\;
  \int_{\Omega_S}\!\Bigl(\int_{\Omega_T} f(\omega_S,\omega_T)\,
      d\mu_T(\omega_T)\Bigr)\,d\mu_S(\omega_S),
\]
where $(\Omega_S,\mu_S):=\bigotimes_{v\in S}(\Omega_v,\mu_v)$ and
$(\Omega_T,\mu_T):=\bigotimes_{v\in T}(\Omega_v,\mu_v)$.
\end{lemma}

\begin{remark}
Lemma~\ref{lem:fubini-product} is an immediate specialization of the
classical Fubini--Tonelli theorem for product measures on finite
products; see, for example, \cite[Section~2.5, Theorems~2.37 and~2.39]{Folland1999}.
\end{remark}

\begin{theorem}[Exponential separation at constant/sublinear depth ]
\label{thm:generic-const-app}
Let $N=n^2$ and let the baseline be as above. Let $p=\alpha_n n$, where
$\alpha_n\in(0,1)$ may depend on $n$. For each row event $E_i$, let $S_i(p)$
denote its depth-$p$ row light cone and define
\[
  \chi_p
  :=
  \max_{v\in[N]}
  \left|\{\,i\in[n]: v\in S_i(p)\,\}\right|.
\]
\begin{enumerate}
\item[(A)] \textbf{1D rows, nearest–row coupling.}
\[
  P_{p}^{\mathrm{(gen)}}
  \;\le\;
  \Bigl[C(2p{+}1)\,2^{-(n-1)}\Bigr]^{n/\chi_p}.
\]
In particular, since $\chi_p\le 2p+1$,
\[
  P_{p}^{\mathrm{(gen)}}
  \;\le\;
  \Bigl[C(2p{+}1)\,2^{-(n-1)}\Bigr]^{n/(2p+1)}.
\]

\item[(B)] \textbf{Bounded intra–row degree $\Delta_{\rm row}$.}
In general,
\[
  P_{p}^{\mathrm{(gen)}}
  \;\le\;
  \Bigl[C\,\Delta_{\rm row}^{\,p}\,2^{-(n-1)}\Bigr]^{n/\chi_p},
  \qquad
  \chi_p\lesssim \min\{n,\Delta_{\rm row}^{\,p}\}.
\]
\end{enumerate}
Hence an exponential CE–vs–generic separation holds whenever
\[
  \frac{n}{\chi_p}
  \left[
    (n-1)\ln 2
    -
    \ln\!\bigl(CW_{\rm row}(p)\bigr)
  \right]
  -
  n\ln n
  \to +\infty.
\]
For 1D nearest–row coupling, this includes
\[
  p=o\!\left(\frac{n}{\ln n}\right).
\]
\end{theorem}


\begin{proof}
Let $E_i$ be the indicator that row $i$ is one–hot, i.e., $E_i:=P_{\mathrm{row},i}$.
By Lemma~\ref{lem:row-valid-general}, for every $i$ we have
\[
  \mathbb{E}[E_i]\le C(2p{+}1)\,2^{-(n-1)}
\]
in case (A), and
\[
  \mathbb{E}[E_i]\le C\,\Delta_{\rm row}^{\,p}\,2^{-(n-1)}
\]
in case (B). Moreover, $E_i$ depends only on the coordinates in its
depth-$p$ row light cone $S_i(p)$. Define
\[
  \chi_p
  :=
  \max_{v\in[N]}
  \left|\{\,i\in[n]: v\in S_i(p)\,\}\right|.
\]
Applying the dependency–graph Hölder bound \eqref{eq:finner} to
$\prod_{i=1}^n E_i$ gives
\[
  P_p^{(\mathrm{gen})}
  \;=\; \mathbb{E}\!\Bigl[\prod_{i=1}^n E_i\Bigr]
  \;\le\; \prod_{i=1}^n \bigl(\mathbb{E}[E_i^{\chi_p}]\bigr)^{1/\chi_p}
  \;=\; \prod_{i=1}^n \bigl(\mathbb{E}[E_i]\bigr)^{1/\chi_p},
\]
where the last equality uses $E_i^{\chi_p}=E_i$. Substituting the per–row
estimates gives
\[
  P_p^{(\mathrm{gen})}
  \;\le\;
  \Bigl[C(2p{+}1)\,2^{-(n-1)}\Bigr]^{n/\chi_p}
\]
in case (A), and
\[
  P_p^{(\mathrm{gen})}
  \;\le\;
  \Bigl[C\,\Delta_{\rm row}^{\,p}\,2^{-(n-1)}\Bigr]^{n/\chi_p}
\]
in case (B). In the 1D nearest–row model,
\[
  \chi_p\le \min\{n,2p+1\},
\]
while in the bounded-degree case,
\[
  \chi_p\lesssim \min\{n,\Delta_{\rm row}^{\,p}\}.
\]
The separation condition follows by comparing this overlap-corrected generic
upper bound with
\[
  P_{1}^{\mathrm{(CE)}}=\Omega(n^{-n})
  =
  \exp[-O(n\ln n)].
\]
Thus the ratio $P_{1}^{\mathrm{(CE)}}/P_{p}^{\mathrm{(gen)}}$ grows
exponentially whenever
\[
  \frac{n}{\chi_p}
  \left[
    (n-1)\ln 2
    -
    \ln\!\bigl(CW_{\rm row}(p)\bigr)
  \right]
  -
  n\ln n
  \to +\infty.
\]
For 1D nearest–row coupling this condition is satisfied whenever
$p=o(n/\ln n)$.
\end{proof}




% From the lightcone analyses, generic QAOA expands correlations at most one step per cost layer in $G(H_C)$. On the other hand, CE--QAOA first \emph{collapses} each visited block to a clique in one shot, then uses the (often dense) penalty couplings to jump across blocks; for permutation penalties this reaches every block in one layer and, after the following mixer, populates entire blocks. %Thus CE-QAOA has a strictly larger correlation radius (in the block graph) and a dramatically larger correlated \emph{volume} at the same depth; it effectively “sees’’ the full problem graph in $O(1)$ depth for collision penalties. See detailed analysis in Appendix \ref{sec:ce-lightcone}




\section{Conclusion and Outlook}
\label{sec:conclusion}

This work reframes the computational power of shallow variational quantum algorithms through the lens of \emph{feasibility concentration}. For permutation encodings on \(n^2\) qubits we prove that the generic QAOA cannot lift feasible mass much beyond the uniform baseline. The ceiling persists at constant depth and, under bounded degree, up to linear depth. In the same regime a minimal constraint–enhanced surrogate, CE–QAOA, achieves a depth-matched and angle-robust exponential gain in feasible mass. The separation follows from three transparent controls on the output distribution—angle averaging with low order moments, Boolean harmonic analysis, and light-cone locality. 

The CE–QAOA kernel is not engineered here. We use it merely as a clean illustration of problem–algorithm co-design as a route to exponential enhancements. As we have demonstrated, even the minimal symmetry based bound we explored here was enough to achieve exponential enhancement. These results suggest that a simple design rule for globally constrained problems is to build feasibility into the ansatz and make the mixer operate \emph{within} the constraint manifold. This converts a global coordination task that shallow generic circuits cannot sustain into a local exploration inside the encoded space, where light-cones grow at the block level and the baseline is information-theoretically favorable. Typical-angle and fourth-moment controls also serve as diagnostics when a generic shallow circuit appears to “fail to train.” The bounds here show that the obstacle is structural, not an artifact of optimizer choice, sampling noise, or unlucky initialization.

Our assumptions are mild but explicit. We work with diagonal commuting costs of arbitrary degree and the standard \(X\) mixer. We analyze permutation encodings because they capture a broad family of assignment and routing objectives and make the obstruction sharp. CE-QAOA–style kernels offer one principled path in this broader space, but other constraint-preserving designs are possible and worth exploring. We conclude that while feasibility concentration is the bottleneck for generic shallow VQA on global constraints, algorithms that \emph{respect} those constraints in part or in full convert that bottleneck into a computational resource and unlock exponential gains. A natural next step is to build the same harmonic and Boolean–cube toolkit inside the CE–QAOA manifold to track how the block–$XY$ mixer reshapes Fourier weight, and identify which degrees are actually reachable at fixed depth. This would give matching upper bounds for CE–QAOA and expose its ultimate limitations rather than only its advantages.

% Beyond this baseline enhancement, we showed a parameter-transfer amplification result that combines the spirit of warmstarting QAOA and parameter transfer in QAOA. If \((\beta^\star,\gamma^\star)\) maximize the feasible mass for generic \(p{=}1\), then reusing those same angles inside the encoded kernel boosts success by at least
% \(
% P^{(\mathrm{CE})}\ \ge\ \sqrt{2\pi n}\,e^{\,n}\;P^{(\mathrm{gen})}_{\max},
% \)
% which makes parameter reuse a principled and testable strategy rather than a heuristic.


\section*{Data Availability Statement} All data in this paper and the Python implementation in Qiskit are made available here \url{https://doi.org/10.5281/zenodo.17701648}.


\section*{Conflict of Interests}
All authors declare no competing interests.

\section*{Author contributions}
C.O. wrote the main paper and prepared all the figures, K.M. provided guidance during the project. All the authors reviewed the paper.

\section*{Funding}
This research received no external funding.






\appendix

\printbibliography
\end{document}



